\chapter[Algebraic Tori and Their Representation Theory]
{Algebraic Tori and Their\\ Representation Theory}
\label{chap:algebraic-tori}

\begin{chapteroverview}
Why are monomials so central in toric geometry?

The answer begins with representation theory. Algebraic tori are among
the simplest algebraic groups, and their finite-dimensional
representations are completely controlled by monomial functions. In
this chapter we develop the basic language of characters,
one-parameter subgroups, weights, and dual lattices. The examples of
\(\Gm\), \(\Gm^2\), and the dense torus of \(\PP^2\) remain in the
foreground throughout.

Most structural results about representations are stated without
proof. References and proofs may be supplied from standard texts on
linear algebraic groups; our goal here is to establish the language
and geometric intuition needed for toric varieties.
\end{chapteroverview}

\begin{threedictionaries}
The central objects of the chapter already come in parallel families.

\medskip

\begin{center}
\small
\begin{tabularx}{\linewidth}{@{}
  >{\raggedright\arraybackslash}X
  >{\raggedright\arraybackslash}X
  >{\raggedright\arraybackslash}X@{}}
\toprule
\textbf{Algebra}
&
\textbf{Geometry}
&
\textbf{Combinatorics}
\\
\midrule
the group algebra \(\Bbbk[M]\)
&
the algebraic torus \(T\)
&
the character lattice \(M\)
\\[1mm]
characters \(T\to\Gm\)
&
monomial functions on \(T\)
&
lattice points \(m\in M\)
\\[1mm]
one-parameter subgroups \(\Gm\to T\)
&
one-parameter motions and limits
&
lattice points \(u\in N\)
\\[1mm]
the composition
\(\character{m}\circ\oneps{u}\)
&
the order of vanishing or pole at \(t=0\)
&
the pairing \(\pair{m}{u}\)
\\
\bottomrule
\end{tabularx}
\end{center}
\end{threedictionaries}

\section{Why ``torus''?}
\label{sec:why-torus}
\index{torus!topological and algebraic}
\index{algebraic torus}

The word \emph{torus} comes to algebraic geometry from topology. A
topologist begins with the circle
\[
S^1=\{z\in\CC:|z|=1\}
\]
and calls a product
\[
(S^1)^n
\]
an \(n\)-dimensional torus. When \(n=2\), this is the familiar surface
of a doughnut; in higher dimensions, the product description is more
useful than any attempted drawing.

What does this have to do with the algebraic group
\[
\Gm=\Spec\Bbbk[x,x^{-1}]?
\]
The connection is easiest to see over \(\CC\).

\subsection{From the circle to the punctured complex line}

The complex points of \(\Gm\) are
\[
\Gm(\CC)=\CC^\times=\CC\setminus\{0\}.
\]
Every nonzero complex number has a polar decomposition
\[
z=r\,e^{i\theta},
\qquad
r>0.
\]
Thus, as a topological space,
\[
\CC^\times\simeq S^1\times\RR_{>0}.
\]
The second factor is contractible. We may continuously push every
point radially toward the unit circle by
\[
H(z,s)
=
\left((1-s)+\frac{s}{|z|}\right)z,
\qquad
0\leq s\leq1.
\]
At \(s=0\), this is the identity; at \(s=1\), it sends
\(z\) to \(z/|z|\in S^1\). In topological language,
\(\CC^\times\) deformation retracts onto \(S^1\).

\begin{figure}[ht]
\centering
\begin{tikzpicture}[scale=1.25]
  \draw[->,black!55] (-2.6,0) -- (2.8,0)
    node[right] {\(\operatorname{Re}\)};
  \draw[->,black!55] (0,-2.35) -- (0,2.55)
    node[above] {\(\operatorname{Im}\)};

  \draw[very thick,PredictionTeal] (0,0) circle (1.35);
  \node[PredictionTeal,above right] at (0.95,0.95) {\(S^1\)};

  \fill[WarningRed] (0,0) circle (2.2pt);
  \node[below right,WarningRed] at (0.05,-0.05) {removed point};

  \foreach \ang/\rad in {18/2.15,55/1.85,125/2.05,205/1.95,310/2.20}{
    \coordinate (P) at ({\rad*cos(\ang)},{\rad*sin(\ang)});
    \coordinate (Q) at ({1.35*cos(\ang)},{1.35*sin(\ang)});
    \fill[DictionaryBlue] (P) circle (1.8pt);
    \draw[thick,-{Latex[length=2mm]},ExampleGold] (P) -- (Q);
  }

  \node[align=center] at (4.65,0.75)
    {\(\displaystyle z\longmapsto \frac{z}{|z|}\)\\[1mm]
     radial contraction};
  \draw[very thick,-{Latex[length=2.5mm]},DictionaryBlue]
    (2.55,0.35) -- (3.55,0.55);
\end{tikzpicture}
\caption{The punctured complex line retracts radially onto the unit
circle. This topological circle is the compact core behind the name
\(\Gm\).}
\label{fig:gm-retracts-to-circle}
\end{figure}

Taking products gives
\[
(\CC^\times)^n
\simeq
(S^1)^n\times(\RR_{>0})^n,
\]
and therefore
\[
(\CC^\times)^n
\quad\text{deformation retracts onto}\quad
(S^1)^n.
\]
This is why \(\Gm^n\) is called an algebraic torus.

\begin{warningbox}
The complex algebraic torus \((\CC^\times)^n\) is not itself compact.
It contains the compact topological torus \((S^1)^n\) as a deformation
retract. The word ``torus'' records this topological relationship; it
does not mean that the algebraic variety is literally a compact
doughnut.
\end{warningbox}

\subsection{The algebraic definition}

For most of this course, the base field \(\Bbbk\) is algebraically
closed.

\begin{definition}
An \(n\)-dimensional algebraic torus over \(\Bbbk\) is an algebraic
group isomorphic to
\[
\Gm^n.
\]
\end{definition}

The word \emph{algebraic group} means that multiplication and inversion
are given by morphisms of algebraic varieties. On points, the group law
of \(\Gm^n\) is coordinatewise multiplication:
\[
(t_1,\ldots,t_n)(s_1,\ldots,s_n)
=
(t_1s_1,\ldots,t_ns_n),
\]
the identity is
\[
(1,\ldots,1),
\]
and inversion is
\[
(t_1,\ldots,t_n)^{-1}
=
(t_1^{-1},\ldots,t_n^{-1}).
\]
The coordinate ring
\[
\Bbbk[t_1^{\pm1},\ldots,t_n^{\pm1}]
\]
is exactly what allows these inverse coordinates to remain regular.

\begin{observation}
Over a field that is not algebraically closed, one distinguishes
between \emph{split} tori, which are already isomorphic to
\(\Gm^n\), and tori that become isomorphic to \(\Gm^n\) only after
extending the base field. The distinction is important in arithmetic
geometry, but it will not play a role in the main development of these
notes.
\end{observation}

\subsection{Examples close to home}

The first examples are already familiar.

\begin{example}
The one-dimensional torus is
\[
\Gm=\AAA^1\setminus\{0\}.
\]
As an affine variety,
\[
\Gm=\Spec\Bbbk[x,x^{-1}].
\]
\end{example}

\begin{example}
The dense open subset of \(\PP^2\) studied in
\cref{sec:p2-dense-torus} is
\[
T
=
\{[x_0:x_1:x_2]:x_0x_1x_2\neq0\}
\simeq\Gm^2.
\]
After normalizing \(x_0=1\), its group coordinates are
\[
\left(\frac{x_1}{x_0},\frac{x_2}{x_0}\right).
\]
Thus Chapter~0 did not merely exhibit an open subset isomorphic to
\(\Gm^2\): it exhibited the algebraic torus that acts on all of
\(\PP^2\).
\end{example}

\begin{example}
The group of invertible diagonal \(n\times n\) matrices is an
\(n\)-dimensional torus:
\[
\left\{
\begin{pmatrix}
t_1&&0\\
&\ddots&\\
0&&t_n
\end{pmatrix}
:
t_i\in\Bbbk^\times
\right\}
\simeq
\Gm^n.
\]
The determinant-one diagonal matrices form an
\((n-1)\)-dimensional torus:
\[
\{(t_1,\ldots,t_n)\in\Gm^n:t_1\cdots t_n=1\}
\simeq
\Gm^{n-1}.
\]
For example, one may solve for \(t_n\) as
\[
t_n=(t_1\cdots t_{n-1})^{-1}.
\]
\end{example}

\begin{pauseandthink}
The affine line with addition,
\[
\Ga=(\AAA^1,+),
\]
is also a one-dimensional algebraic group. Could it be isomorphic to
an algebraic torus? Try to approach the question both algebraically
and, over \(\CC\), topologically.
\end{pauseandthink}

The algebraic answer is that \(\Ga\) is not isomorphic to \(\Gm\):
its coordinate ring is \(\Bbbk[x]\), not
\(\Bbbk[x,x^{-1}]\), and its group law is addition rather than
multiplication. Over \(\CC\), the spaces are topologically different
as well:
\[
\CC\quad\text{is contractible,}
\qquad
\CC^\times\quad\text{retracts onto }S^1.
\]

\subsection{Why monomials are waiting for us}

The multiplication law on a torus strongly favors monomials. On
\(\Gm\), the functions
\[
t\longmapsto t^m,
\qquad
m\in\ZZ,
\]
satisfy
\[
(ts)^m=t^ms^m.
\]
On \(\Gm^n\), the Laurent monomials
\[
t_1^{m_1}\cdots t_n^{m_n}
\]
have the same property:
\[
(t_1s_1)^{m_1}\cdots(t_ns_n)^{m_n}
=
\bigl(t_1^{m_1}\cdots t_n^{m_n}\bigr)
\bigl(s_1^{m_1}\cdots s_n^{m_n}\bigr).
\]

\begin{prediction}
Suppose
\[
f:\Gm^n\longrightarrow\Gm
\]
is both a regular map and a group homomorphism. What forms can such
maps take? Explore the cases \(n=1\) and \(n=2\) first, and try to
organize the examples you find.
\end{prediction}

The answer is yes. Such maps are called \emph{characters}, and their
exponent vectors form a lattice. This simple fact is the entrance to
the representation theory of algebraic tori and the explanation for
the persistent appearance of monomials in Chapter~0.


\section{The algebraic torus
\texorpdfstring{\(\Gm^n\)}{Gm n}}
\label{sec:algebraic-torus-gn}
\index{algebraic torus}
\index{Laurent polynomial}
\index{group algebra}

The notation
\[
\Gm^n
\]
is compact, but it hides several compatible structures. It is at once

\[
\text{an affine variety,}
\qquad
\text{an algebraic group,}
\qquad
\text{and a space carrying many monomial functions.}
\]

In this section we put these structures on the table. The point is not
to introduce a large amount of machinery. It is to understand exactly
why Laurent polynomials are the natural algebraic language of a torus.

\subsection{Removing the coordinate hyperplanes}

Begin with affine space
\[
\AAA^n=\Spec\Bbbk[x_1,\ldots,x_n].
\]
The algebraic torus is the open subset on which no coordinate vanishes:
\[
\Gm^n
=
\{(t_1,\ldots,t_n)\in\AAA^n:t_1\cdots t_n\neq0\}.
\]
Equivalently,
\[
\Gm^n
=
D(x_1\cdots x_n)
=
\AAA^n\setminus V(x_1\cdots x_n).
\]
Since
\[
V(x_1\cdots x_n)
=
V(x_1)\cup\cdots\cup V(x_n),
\]
we obtain \(\Gm^n\) by removing the \(n\) coordinate hyperplanes.

For \(n=2\), this is the familiar picture
\[
\Gm^2
=
\AAA^2\setminus\bigl(V(x)\cup V(y)\bigr).
\]

\begin{figure}[ht]
\centering
\begin{tikzpicture}[scale=1.05]
  \fill[PastelYellow] (-2.7,-2.3) rectangle (2.7,2.3);

  \draw[very thick,WarningRed] (-2.7,0) -- (2.7,0);
  \draw[very thick,DictionaryBlue] (0,-2.3) -- (0,2.3);

  \node[align=center] at (1.35,1.25)
    {\(\displaystyle \Gm^2\)\\
     \(x\neq0,\ y\neq0\)};

  \node[below,WarningRed] at (1.75,-0.08)
    {\(V(y)\)};
  \node[rotate=90,above,DictionaryBlue] at (-0.08,1.55)
    {\(V(x)\)};

  \fill[black] (0,0) circle (1.8pt);
\end{tikzpicture}
\caption{The torus \(\Gm^2\) is obtained from \(\AAA^2\) by removing
the two coordinate axes. The colored lines are not part of the torus.}
\label{fig:gm2-complement-coordinate-axes}
\end{figure}

This description immediately explains why the coordinates \(x_i\)
become invertible on the torus. None of them vanishes there, so each
reciprocal
\[
x_i^{-1}
\]
is a regular function on \(\Gm^n\).

\subsection{Laurent polynomials are the regular functions}

The coordinate ring of \(\Gm^n\) is obtained by inverting the product
\(x_1\cdots x_n\):
\[
\Bbbk[\Gm^n]
=
\Bbbk[x_1,\ldots,x_n]_{x_1\cdots x_n}.
\]
Since inverting the product makes every factor invertible, this ring is
more familiarly written as
\[
\Bbbk[\Gm^n]
=
\Bbbk[x_1^{\pm1},\ldots,x_n^{\pm1}].
\]

Its elements are Laurent polynomials:
\[
f
=
\sum_{m\in F} c_m x_1^{m_1}\cdots x_n^{m_n},
\]
where

\[
F\subseteq\ZZ^n
\]
is finite and the exponents \(m_i\) may be positive, zero, or negative.

\begin{example}
The function
\[
f(x,y)=x^{-2}y^3+4xy^{-1}-7
\]
is regular on \(\Gm^2\). It is not regular on all of \(\AAA^2\),
because its negative exponents produce poles along the coordinate
axes.
\end{example}

Thus passing from affine space to the torus has one simple algebraic
effect:
\[
\begin{array}{c@{\qquad\rightsquigarrow\qquad}c}
\AAA^n & \Gm^n\\[1mm]
\Bbbk[x_1,\ldots,x_n]
&
\Bbbk[x_1^{\pm1},\ldots,x_n^{\pm1}].
\end{array}
\]
Removing the coordinate hyperplanes allows negative exponents.

\begin{pauseandthink}
Consider the Laurent monomial
\[
x_1^{-3}x_2^2.
\]
On which open subset of \(\AAA^2\) is it regular? What is the smallest
standard principal open subset on which it becomes a regular function?
\end{pauseandthink}

Only the negative exponent creates a denominator. Hence the monomial
is regular on
\[
D(x_1)=\{x_1\neq0\}.
\]
There is no need to remove \(V(x_2)\), because the positive power
\(x_2^2\) is regular everywhere.

\subsection{The group law is algebraic}

On points, multiplication is coordinatewise:
\[
\mu:\Gm^n\times\Gm^n\longrightarrow\Gm^n,
\]
\[
\mu\bigl((t_1,\ldots,t_n),(s_1,\ldots,s_n)\bigr)
=
(t_1s_1,\ldots,t_ns_n).
\]
The inverse map is
\[
\iota:\Gm^n\longrightarrow\Gm^n,
\qquad
(t_1,\ldots,t_n)
\longmapsto
(t_1^{-1},\ldots,t_n^{-1}),
\]
and the identity element is
\[
e=(1,\ldots,1).
\]

\begin{pauseandthink}
Why are multiplication and inversion regular maps on \(\Gm^n\)?
Write each output coordinate as a regular function on the source, and
identify where the fact that the coordinates \(t_i\) are invertible is
being used.
\end{pauseandthink}

Multiplication and inversion are regular because their coordinate
functions are Laurent monomials.

\begin{figure}[ht]
\centering
\begin{tikzpicture}[
  operation/.style={
    rounded corners,
    draw=black!55,
    minimum width=3.45cm,
    minimum height=1.35cm,
    align=center
  }
]
  \node[operation,fill=PastelYellow] (mult) at (-3.8,0)
    {\textbf{Multiplication}\\[1mm]
     \(\displaystyle
     (t_i),(s_i)\longmapsto(t_is_i)\)};

  \node[operation,fill=DictionaryBlue!8] (inv) at (0,0)
    {\textbf{Inversion}\\[1mm]
     \(\displaystyle
     (t_i)\longmapsto(t_i^{-1})\)};

  \node[operation,fill=SurfaceGreen!9] (id) at (3.8,0)
    {\textbf{Identity}\\[1mm]
     \(\displaystyle
     e=(1,\ldots,1)\)};

  \node[below=8mm,align=center] at (0,0)
    {all formulas are built from Laurent monomials};
\end{tikzpicture}
\caption{The algebraic group structure of \(\Gm^n\). Laurent monomials
make both multiplication and inversion regular.}
\label{fig:gm-group-operations}
\end{figure}

There is an equivalent description on coordinate rings. A regular map
of affine varieties produces a homomorphism of coordinate rings in the
opposite direction. For multiplication, inversion, and the identity,
the coordinate functions satisfy
\[
\begin{aligned}
\mu^*(x_i)&=x_i\otimes x_i,\\
\iota^*(x_i)&=x_i^{-1},\\
e^*(x_i)&=1.
\end{aligned}
\]
The first formula says that the \(i\)-th coordinate of a product is
the product of the two \(i\)-th coordinates.

\begin{observation}
The three displayed formulas are the beginning of the Hopf algebra
description of an algebraic group. We will not need that language in
the main development. What matters here is the concrete principle:
the group law on points becomes a rule for pulling back regular
functions.
\end{observation}

\subsection{Maps into \texorpdfstring{\(\Gm\)}{Gm} and invertible functions}

There is a useful general principle. Let \(X\) be an affine variety.
A morphism
\[
f:X\longrightarrow\Gm
\]
is the same thing as an invertible regular function on \(X\).

Indeed, the coordinate \(z\) on \(\Gm\) is invertible, so its pullback
\[
f^*(z)
\]
must be a unit in \(\Bbbk[X]\). Conversely, a unit
\[
g\in\Bbbk[X]^\times
\]
defines a morphism
\[
X\longrightarrow\Gm,
\qquad
p\longmapsto g(p).
\]

Specializing to \(X=\Gm^n\), maps
\[
\Gm^n\longrightarrow\Gm
\]
are therefore controlled by the units of the Laurent polynomial ring
\[
\Bbbk[x_1^{\pm1},\ldots,x_n^{\pm1}].
\]
In the next section, we impose one additional condition: the map should
respect multiplication. This will isolate the characters of the torus.

\begin{prediction}
Among all invertible Laurent polynomials, which ones should define
group homomorphisms
\[
\Gm^n\longrightarrow\Gm?
\]
Try testing sums of monomials, scalar multiples of monomials, and pure
Laurent monomials against the identity
\[
f(ts)=f(t)f(s).
\]
\end{prediction}

\subsection{The dense torus of \texorpdfstring{\(\PP^2\)}{P2} revisited}

The dense torus of \(\PP^2\) is
\[
T
=
\{[x_0:x_1:x_2]:x_0x_1x_2\neq0\}.
\]
On the chart \(U_0\), it has coordinates
\[
z_1=\frac{x_1}{x_0},
\qquad
z_2=\frac{x_2}{x_0},
\]
and hence
\[
\Bbbk[T]
=
\Bbbk[z_1^{\pm1},z_2^{\pm1}].
\]

The other projective charts give different Laurent coordinates on the
same torus. On \(U_1\), for instance,
\[
w_0=\frac{x_0}{x_1}=z_1^{-1},
\qquad
w_2=\frac{x_2}{x_1}=z_1^{-1}z_2.
\]
Thus the change of coordinates is monomial:
\[
(z_1,z_2)
\longmapsto
(z_1^{-1},z_1^{-1}z_2).
\]

This is our first indication that the torus itself does not come with a
preferred coordinate system. Different choices of coordinates are
related by invertible monomial transformations. Later we will encode
such changes by matrices in
\[
\operatorname{GL}(n,\ZZ).
\]

\begin{threedictionaries}
The algebraic torus may be read in three compatible ways.

\medskip

\noindent
\begin{minipage}[t]{0.30\linewidth}
\centering
\textbf{Algebra}\par
\smallskip
Its coordinate ring is
\[
\Bbbk[x_1^{\pm1},\ldots,x_n^{\pm1}].
\]
Regular functions are finite sums of Laurent monomials.
\end{minipage}
\hfill
\begin{minipage}[t]{0.30\linewidth}
\centering
\textbf{Geometry}\par
\smallskip
It is the complement of the coordinate hyperplanes:
\[
\Gm^n
=
\AAA^n\setminus V(x_1\cdots x_n).
\]
\end{minipage}
\hfill
\begin{minipage}[t]{0.30\linewidth}
\centering
\textbf{Combinatorics}\par
\smallskip
Laurent monomials are indexed by exponent vectors
\[
m=(m_1,\ldots,m_n)\in\ZZ^n.
\]
\end{minipage}
\end{threedictionaries}

\begin{guidingquestion}
Regular functions on a torus are sums of Laurent monomials. Which of
these functions are compatible with the group law?
\end{guidingquestion}


\section{Monomials and characters}
\label{sec:monomials-characters}
\index{character}
\index{monomial!as character}
\index{representation!one-dimensional}

At the end of the previous section, we asked which regular maps
\[
f:\Gm^n\longrightarrow\Gm
\]
are compatible with multiplication. There are many regular functions
on a torus, but the group law is highly restrictive. It singles out
the Laurent monomials.

\subsection{Multiplicative regular functions}

\begin{definition}
A \emph{character} of an algebraic torus \(T\) is a morphism of
algebraic groups
\[
\chi:T\longrightarrow\Gm.
\]
Thus \(\chi\) is a regular map satisfying
\[
\chi(ts)=\chi(t)\chi(s)
\qquad
\text{and}
\qquad
\chi(e)=1.
\]
\end{definition}

For \(T=\Gm^n\), every exponent vector
\[
m=(m_1,\ldots,m_n)\in\ZZ^n
\]
defines a character
\[
\character{m}(t_1,\ldots,t_n)
=
t_1^{m_1}\cdots t_n^{m_n}.
\]
It will often be convenient to abbreviate this Laurent monomial as
\[
t^m.
\]

Negative exponents cause no problem: every coordinate \(t_i\) is
invertible on the torus. Moreover,
\[
\character{m}(ts)
=
(t_1s_1)^{m_1}\cdots(t_ns_n)^{m_n}
=
\character{m}(t)\character{m}(s),
\]
so \(\character{m}\) is indeed a group homomorphism.

\begin{example}
On \(\Gm^2\), the exponent vector
\[
m=(2,-3)
\]
defines the character
\[
\character{(2,-3)}(t_1,t_2)
=
t_1^2t_2^{-3}.
\]
The vector \(m=(0,0)\) gives the constant character
\[
\character{0}=1.
\]
\end{example}

\begin{figure}[ht]
\centering
\begin{tikzpicture}[
  mainbox/.style={
    rounded corners,
    draw=black!55,
    text width=3.0cm,
    minimum height=1.45cm,
    align=center,
    inner sep=5pt
  }
]
  \node[mainbox,fill=DictionaryBlue!8] (lattice) at (-4.0,0)
    {\(\displaystyle m=(m_1,\ldots,m_n)\)\\[1mm]
     an integral exponent vector};

  \node[mainbox,fill=PastelYellow] (monomial) at (0,0)
    {\(\displaystyle t^m
      =t_1^{m_1}\cdots t_n^{m_n}\)\\[1mm]
     a Laurent monomial};

  \node[mainbox,fill=SurfaceGreen!9] (character) at (4.0,0)
    {\(\displaystyle \character{m}:\Gm^n\to\Gm\)\\[1mm]
     a character};

  \draw[very thick,-{Latex[length=2.4mm]},DictionaryBlue]
    (lattice.east) -- (monomial.west);
  \draw[very thick,-{Latex[length=2.4mm]},PredictionTeal]
    (monomial.east) -- (character.west);
\end{tikzpicture}
\caption{For a split torus, an integral exponent vector, a Laurent
monomial, and a character are three descriptions of the same object.}
\label{fig:exponent-monomial-character}
\end{figure}

\subsection{There are no other characters}

The examples above give every character of \(\Gm^n\).

\begin{theorem}[Characters of a split torus]
\label{thm:characters-split-torus}
Every character
\[
\chi:\Gm^n\longrightarrow\Gm
\]
has the form
\[
\chi(t_1,\ldots,t_n)
=
t_1^{m_1}\cdots t_n^{m_n}
\]
for a unique vector
\[
m=(m_1,\ldots,m_n)\in\ZZ^n.
\]
\end{theorem}

The statement follows from two elementary facts.

First, a morphism
\[
\Gm^n\longrightarrow\Gm
\]
is an invertible regular function on \(\Gm^n\). The units in the
Laurent polynomial ring are precisely
\[
c\,t_1^{m_1}\cdots t_n^{m_n},
\qquad
c\in\Bbbk^\times,\quad m_i\in\ZZ.
\]
Second, a group homomorphism must send the identity to the identity.
Evaluating at
\[
e=(1,\ldots,1)
\]
forces
\[
c=1.
\]
Thus the multiplicative regular functions are exactly the pure Laurent
monomials.

\begin{observation}
The description of the units in a Laurent polynomial ring is standard.
One may prove it using the \(\ZZ^n\)-grading, or by comparing the
extreme exponent vectors in a product. We will use the result freely;
the Notes and references section lists standard sources for this
background.
\end{observation}

\begin{pauseandthink}
Why does a scalar multiple
\[
c\,t^m
\qquad(c\in\Bbbk^\times)
\]
fail to be a character unless \(c=1\)? Can you see the obstruction
both by evaluating at the identity and by testing the equation
\[
f(ts)=f(t)f(s)?
\]
\end{pauseandthink}

Not every regular function is a character. For instance,
\[
t_1+t_2
\]
is a regular function on \(\Gm^2\), but it can vanish and therefore
does not even define a map to \(\Gm\). A sum of several Laurent
monomials usually loses the multiplicative behavior that characters
require.

\subsection{What the exponents record}

Once coordinates have been chosen, the torus sits naturally inside
affine space:
\[
\Gm^n
=
\AAA^n\setminus\bigl(H_1\cup\cdots\cup H_n\bigr),
\qquad
H_i=V(x_i).
\]
A character
\[
\character{m}
=
x_1^{m_1}\cdots x_n^{m_n}
\]
is a regular function on \(\Gm^n\), but it may have zeros or poles
when regarded as a rational function on the surrounding affine space
\(\AAA^n\).

The exponent \(m_i\) records exactly what happens along the coordinate
hyperplane \(H_i\):
\[
\ord_{H_i}(\character{m})=m_i.
\]
Thus
\[
\begin{array}{c|c}
m_i>0
&
\character{m}\text{ vanishes to order }m_i\text{ along }H_i,
\\
m_i=0
&
\character{m}\text{ has neither a zero nor a pole along }H_i,
\\
m_i<0
&
\character{m}\text{ has a pole of order }-m_i\text{ along }H_i.
\end{array}
\]
Equivalently, as a rational function on \(\AAA^n\),
\[
\operatorname{div}_{\AAA^n}(\character{m})
=
m_1H_1+\cdots+m_nH_n.
\]

This interpretation depends on the chosen coordinates: the affine
space \(\AAA^n\) and its coordinate hyperplanes arise from the chosen
identification \(T\simeq\Gm^n\). Later, when we work intrinsically,
orders obtained by composing characters with one-parameter subgroups
will be encoded by the pairing between the character and
one-parameter subgroup lattices.

\begin{example}
Let
\[
\chi(t_1,t_2,t_3)=t_1^2t_2^{-1}t_3^4.
\]
As a rational function on \(\AAA^3\),
\[
\operatorname{div}_{\AAA^3}(\chi)
=
2H_1-H_2+4H_3.
\]
Thus \(\chi\) vanishes to order \(2\) along \(H_1\), has a simple pole
along \(H_2\), and vanishes to order \(4\) along \(H_3\).
\end{example}

\begin{observation}
For the coordinate one-parameter map
\[
\epsilon_i(s)
=
(1,\ldots,1,s,1,\ldots,1),
\]
one also has
\[
\character{m}\bigl(\epsilon_i(s)\bigr)=s^{m_i}.
\]
The same integer \(m_i\) therefore appears as an order at \(s=0\).
This is a special coordinate case of the character--one-parameter
subgroup pairing developed later in the chapter.
\end{observation}

\begin{prediction}
Suppose no coordinates have been chosen on an abstract torus \(T\).
How might one record all of its characters without first identifying
\(T\) with \(\Gm^n\)? What algebraic structure should the collection
of characters carry?
\end{prediction}

\subsection{Multiplying characters adds exponents}

Characters can be multiplied pointwise:
\[
(\chi\psi)(t)=\chi(t)\psi(t).
\]
The product is again a character. For monomial characters,
\[
\character{m}\character{m'}
=
\character{m+m'}.
\]
Indeed,
\[
t^m t^{m'}
=
t^{m+m'}.
\]
Similarly,
\[
(\character{m})^{-1}
=
\character{-m},
\]
and the identity element is the constant character
\[
\character{0}=1.
\]

Consequently, the characters of \(\Gm^n\) form an abelian group under
multiplication, and the exponent correspondence gives an isomorphism
\[
\ZZ^n
\longrightarrow
\Hom_{\mathrm{alg.grp}}(\Gm^n,\Gm),
\qquad
m\longmapsto\character{m}.
\]
Addition on the left corresponds to multiplication on the right.

\begin{figure}[ht]
\centering
\begin{tikzpicture}[
  formula/.style={
    rounded corners,
    draw=black!55,
    minimum width=3.25cm,
    minimum height=1.2cm,
    align=center
  }
]
  \node[formula,fill=DictionaryBlue!8] (add) at (-3.6,0)
    {\(\displaystyle m+m'\)\\
     addition of exponent vectors};

  \node[formula,fill=PastelYellow] (mult) at (3.6,0)
    {\(\displaystyle
      \character{m}\character{m'}\)\\
     multiplication of characters};

  \draw[very thick,<->,PredictionTeal]
    (add) -- node[above,fill=white,inner sep=2pt]
    {\(\displaystyle
      \character{m+m'}=\character{m}\character{m'}\)}
    (mult);
\end{tikzpicture}
\caption{The group law on characters becomes ordinary addition of
exponent vectors.}
\label{fig:add-exponents-multiply-characters}
\end{figure}

This group is traditionally denoted
\[
X^*(T)
=
\Hom_{\mathrm{alg.grp}}(T,\Gm).
\]
For \(T=\Gm^n\), the chosen coordinates identify it with \(\ZZ^n\).
The next section studies this group intrinsically and explains why it
should be regarded as a lattice.

\subsection{Characters as one-dimensional representations}

A character also determines an action of \(T\) on a one-dimensional
vector space. If \(V=\Bbbk v\), define
\[
t\cdot v=\chi(t)v.
\]
The group law for the action is exactly the identity
\[
\chi(ts)=\chi(t)\chi(s).
\]
Conversely, every algebraic action of \(T\) on a one-dimensional
vector space arises from a character.

\begin{observation}
This is the representation-theoretic reason for the word
\emph{character}: characters are precisely the one-dimensional
algebraic representations of \(T\). This use of the word differs from
the character of a finite-dimensional representation defined as its
trace. Both notions will coexist comfortably once weight-space
decompositions enter the picture.
\end{observation}

\begin{example}
Let \(T=\Gm^2\) act on the one-dimensional vector space
\(V=\Bbbk v\) by
\[
(t_1,t_2)\cdot v=t_1^3t_2^{-2}v.
\]
Then \(V\) is a one-dimensional representation with character
\[
\character{(3,-2)}.
\]
\end{example}

\begin{threedictionaries}
A character is one object seen from three viewpoints.

\medskip

\noindent
\begin{minipage}[t]{0.30\linewidth}
\centering
\textbf{Algebra}\par
\smallskip
It is a group homomorphism
\[
\chi:T\to\Gm
\]
or, equivalently, a one-dimensional representation of \(T\).
\end{minipage}
\hfill
\begin{minipage}[t]{0.30\linewidth}
\centering
\textbf{Geometry}\par
\smallskip
It is a nowhere-vanishing regular function compatible with
multiplication on the torus.
\end{minipage}
\hfill
\begin{minipage}[t]{0.30\linewidth}
\centering
\textbf{Combinatorics}\par
\smallskip
After choosing torus coordinates, it is indexed by an exponent vector
\[
m\in\ZZ^n.
\]
\end{minipage}
\end{threedictionaries}

\begin{guidingquestion}
The identification \(X^*(\Gm^n)\simeq\ZZ^n\) used a chosen coordinate
system. What remains of this description when the torus is presented
without coordinates?
\end{guidingquestion}


\section{The character lattice}
\label{sec:character-lattice}
\index{character lattice}
\index{lattice!character}
\index{group algebra}

For the coordinatized torus \(\Gm^n\), characters were indexed by
vectors in \(\ZZ^n\). The coordinates made this description easy, but
they also hid what was intrinsic. An abstract torus does not come with
preferred functions
\[
t_1,\ldots,t_n.
\]
What it does come with is the collection of all its characters.

\subsection{The intrinsic definition}

\begin{definition}
Let \(T\) be an algebraic torus. Its \emph{character group} is
\[
X^*(T)
=
\Hom_{\mathrm{alg.grp}}(T,\Gm).
\]
The group law is pointwise multiplication:
\[
(\chi\psi)(t)=\chi(t)\psi(t).
\]
\end{definition}

The identity element is the constant character \(1\), and the inverse
of \(\chi\) is the character
\[
\chi^{-1}:t\longmapsto\chi(t)^{-1}.
\]
Thus \(X^*(T)\) is an abelian group.

For \(T=\Gm^n\), the classification of characters from
\cref{thm:characters-split-torus} gives
\[
X^*(\Gm^n)\simeq\ZZ^n.
\]
Under this identification,
\[
m+m'
\quad\longleftrightarrow\quad
\character{m}\character{m'}.
\]

\begin{definition}
A \emph{lattice} is a free abelian group of finite rank.
\end{definition}

\begin{theorem}
If \(T\) is an \(n\)-dimensional algebraic torus, then
\[
M:=X^*(T)
\]
is a lattice of rank \(n\).
\end{theorem}

No preferred isomorphism
\[
M\simeq\ZZ^n
\]
is built into the torus. Such an isomorphism appears only after
coordinates have been chosen.

\begin{observation}
Because our tori are split by definition, the theorem follows by
choosing an isomorphism \(T\simeq\Gm^n\). The important point is not
the existence of such an isomorphism, but the fact that the resulting
lattice \(X^*(T)\) is intrinsic even though a basis is not.
\end{observation}

\subsection{Coordinates are bases of the character lattice}

Suppose
\[
\varphi:T\overset{\sim}{\longrightarrow}\Gm^n
\]
is a choice of torus coordinates. Pulling back the coordinate
characters on \(\Gm^n\) gives characters
\[
\chi_1,\ldots,\chi_n\in X^*(T).
\]
Every character of \(T\) can then be written uniquely as
\[
\chi_1^{m_1}\cdots\chi_n^{m_n},
\qquad
(m_1,\ldots,m_n)\in\ZZ^n.
\]
Thus \(\chi_1,\ldots,\chi_n\) form a basis of the lattice \(M\).

Conversely, a basis
\[
\chi_1,\ldots,\chi_n
\]
of \(M\) defines a map
\[
T\longrightarrow\Gm^n,
\qquad
t\longmapsto
\bigl(\chi_1(t),\ldots,\chi_n(t)\bigr),
\]
and this map is an isomorphism of algebraic tori.

Therefore:
\[
\boxed{
\text{choosing coordinates on }T
\quad\Longleftrightarrow\quad
\text{choosing a basis of }X^*(T).
}
\]

\begin{figure}[ht]
\centering
\begin{tikzpicture}[
  mainbox/.style={
    rounded corners,
    draw=black!55,
    text width=3.25cm,
    minimum height=1.45cm,
    align=center,
    inner sep=5pt
  }
]
  \node[mainbox,fill=PastelYellow] (torus) at (-4.0,0)
    {\textbf{Geometry}\\[1mm]
     a coordinate system\\
     \(T\simeq\Gm^n\)};

  \node[mainbox,fill=DictionaryBlue!8] (basis) at (0,0)
    {\textbf{Algebra}\\[1mm]
     a basis\\
     \(\chi_1,\ldots,\chi_n\) of \(X^*(T)\)};

  \node[mainbox,fill=SurfaceGreen!9] (vectors) at (4.0,0)
    {\textbf{Combinatorics}\\[1mm]
     an identification\\
     \(M\simeq\ZZ^n\)};

  \draw[very thick,<->,ExampleGold]
    (torus.east) -- (basis.west);
  \draw[very thick,<->,PredictionTeal]
    (basis.east) -- (vectors.west);
\end{tikzpicture}
\caption{Choosing coordinates on a torus is the same as choosing a
basis of its character lattice. The lattice is intrinsic; the basis is
not.}
\label{fig:coordinates-basis-character-lattice}
\end{figure}

\begin{pauseandthink}
A vector space does not come with a preferred basis. In what sense is
the same true of a torus? Which formulas from the previous sections
depended on a chosen basis of \(X^*(T)\), and which statements did
not?
\end{pauseandthink}

\subsection{The coordinate ring as a group algebra}

Let \(M\) be a lattice. Its \emph{group algebra} over \(\Bbbk\) is
\[
\Bbbk[M]
=
\bigoplus_{m\in M}\Bbbk\,\character{m},
\]
where the symbols \(\character{m}\) multiply according to
\[
\character{m}\character{m'}
=
\character{m+m'}.
\]
An element of \(\Bbbk[M]\) is a finite sum
\[
\sum_{m\in F} c_m\character{m},
\qquad
F\subseteq M\text{ finite}.
\]

For \(M=\ZZ^n\), choosing the standard basis identifies
\[
\Bbbk[M]
\simeq
\Bbbk[x_1^{\pm1},\ldots,x_n^{\pm1}].
\]
The basis element \(\character{(m_1,\ldots,m_n)}\) becomes the Laurent
monomial
\[
x_1^{m_1}\cdots x_n^{m_n}.
\]

For an arbitrary torus \(T\), evaluation of characters gives a
canonical isomorphism
\[
\Bbbk[M]
\overset{\sim}{\longrightarrow}
\Bbbk[T],
\qquad
\character{m}\longmapsto
\bigl(t\mapsto\character{m}(t)\bigr).
\]
Thus the coordinate ring of a torus is canonically the group algebra
of its character lattice:
\[
\boxed{\Bbbk[T]\simeq\Bbbk[X^*(T)].}
\]

The word \emph{canonically} matters. The Laurent polynomial
presentation
\[
\Bbbk[T]
\simeq
\Bbbk[x_1^{\pm1},\ldots,x_n^{\pm1}]
\]
depends on a basis of \(M\), but the group algebra
\[
\Bbbk[M]
\]
does not.

\begin{example}
Let \(T\subseteq\PP^2\) be the dense torus. On the chart \(U_0\), the
coordinates
\[
z_1=\frac{x_1}{x_0},
\qquad
z_2=\frac{x_2}{x_0}
\]
give a basis \(m_1,m_2\) of \(M=X^*(T)\).

On \(U_1\), the torus coordinates are
\[
w_0=\frac{x_0}{x_1}=z_1^{-1},
\qquad
w_2=\frac{x_2}{x_1}=z_1^{-1}z_2.
\]
In the basis \(m_1,m_2\), these new coordinate characters have exponent
vectors
\[
(-1,0)
\qquad\text{and}\qquad
(-1,1).
\]
They form another basis of the same lattice \(M\).

Likewise, the coordinates on \(U_2\),
\[
r_0=\frac{x_0}{x_2}=z_2^{-1},
\qquad
r_1=\frac{x_1}{x_2}=z_1z_2^{-1},
\]
correspond to the basis vectors
\[
(0,-1)
\qquad\text{and}\qquad
(1,-1).
\]
The torus has not changed; only the chosen basis of its character
lattice has changed.
\end{example}

\subsection{Morphisms reverse the direction on character lattices}

Let
\[
\phi:T_1\longrightarrow T_2
\]
be a morphism of algebraic tori. Every character
\[
\chi:T_2\longrightarrow\Gm
\]
can be pulled back along \(\phi\):
\[
\phi^*(\chi)
=
\chi\circ\phi:T_1\longrightarrow\Gm.
\]
Thus \(\phi\) induces a homomorphism of lattices
\[
\phi^*:X^*(T_2)\longrightarrow X^*(T_1).
\]

Notice the reversal of direction:
\[
T_1\longrightarrow T_2
\qquad\text{produces}\qquad
X^*(T_2)\longrightarrow X^*(T_1).
\]
The character lattice is therefore a contravariant construction.

\begin{example}
Consider the diagonal embedding
\[
\delta:\Gm\longrightarrow\Gm^2,
\qquad
t\longmapsto(t,t).
\]
A character of \(\Gm^2\) has the form
\[
\character{(a,b)}(t_1,t_2)=t_1^at_2^b.
\]
Pulling it back gives
\[
\delta^*(\character{(a,b)})(t)
=
t^{a+b}.
\]
On lattices, the map is therefore
\[
\delta^*:\ZZ^2\longrightarrow\ZZ,
\qquad
(a,b)\longmapsto a+b.
\]
\end{example}

\begin{pauseandthink}
Why must the direction reverse? Start with a character on the target
torus and ask where it can be evaluated after a point of the source
torus has been mapped forward.
\end{pauseandthink}

\subsection{Passing to real vector spaces}

Once the lattice \(M\) has been identified, we form the associated
real vector space
\[
M_{\RR}
=
M\otimes_{\ZZ}\RR.
\]
If \(\rank M=n\), then
\[
M_{\RR}\simeq\RR^n,
\]
but again only after choosing a basis.

The lattice points
\[
M\subseteq M_{\RR}
\]
index characters. Convex combinations of these points do not usually
correspond to characters, but they allow us to build polytopes and
polyhedra. This is why the character lattice is the natural home for
the triangles and polytopes that appeared in Chapter~0.

\begin{guidingquestion}
The character lattice records maps \(T\to\Gm\). What should record maps
in the opposite direction,
\[
\Gm\longrightarrow T?
\]
How might such maps interact with characters?
\end{guidingquestion}

\begin{threedictionaries}
The character lattice packages the intrinsic algebra of a torus.

\medskip

\noindent
\begin{minipage}[t]{0.30\linewidth}
\centering
\textbf{Algebra}\par
\smallskip
\[
\Bbbk[T]\simeq\Bbbk[M].
\]
Characters form a multiplicative basis of the coordinate ring.
\end{minipage}
\hfill
\begin{minipage}[t]{0.30\linewidth}
\centering
\textbf{Geometry}\par
\smallskip
A basis of \(M\) is a coordinate system
\[
T\simeq\Gm^n.
\]
Different coordinates describe the same torus.
\end{minipage}
\hfill
\begin{minipage}[t]{0.30\linewidth}
\centering
\textbf{Combinatorics}\par
\smallskip
\(M\) is a lattice, and
\[
M_{\RR}=M\otimes_{\ZZ}\RR
\]
is the space in which exponent polytopes live.
\end{minipage}
\end{threedictionaries}



\section{One-parameter subgroups}
\label{sec:one-parameter-subgroups}
\index{one-parameter subgroup}
\index{cocharacter lattice}
\index{one-parameter subgroup lattice}
\index{primitive lattice vector}

Characters are morphisms
\[
T\longrightarrow\Gm.
\]
We now reverse the arrow and study morphisms
\[
\Gm\longrightarrow T.
\]
Geometrically, these maps are parameterized algebraic curves inside
the torus: the parameter \(t\) ranges through \(\Gm\). Because the
curve respects the group law, it is natural to think of it as an
algebraic motion. When the torus sits inside a larger variety,
following such a motion toward \(t=0\) often reveals a boundary point.

\subsection{Turning the arrow around}

\begin{definition}
A \emph{one-parameter subgroup} of an algebraic torus \(T\) is a
morphism of algebraic groups
\[
\lambda:\Gm\longrightarrow T.
\]
The collection of all one-parameter subgroups is denoted
\[
X_*(T)
=
\Hom_{\mathrm{alg.grp}}(\Gm,T).
\]
It is also called the \emph{cocharacter group} of \(T\).
\end{definition}

A one-parameter subgroup is therefore a very special parameterized
curve in \(T\). It passes through the identity at \(t=1\), and its
parameterization is compatible with multiplication:
\[
\lambda(tt')=\lambda(t)\lambda(t').
\]
Not every algebraic curve in \(T\) has this property.

For the coordinatized torus \(T=\Gm^n\), every vector
\[
u=(u_1,\ldots,u_n)\in\ZZ^n
\]
defines a one-parameter subgroup
\[
\oneps{u}(t)
=
(t^{u_1},\ldots,t^{u_n}).
\]
Negative exponents are allowed because \(t\neq0\) on \(\Gm\).

\begin{example}
The vector
\[
u=(2,-1,3)
\]
defines
\[
\oneps{u}(t)
=
(t^2,t^{-1},t^3)
\]
in \(\Gm^3\).
\end{example}

The converse also holds.

\begin{theorem}[One-parameter subgroups of a split torus]
\label{thm:one-parameter-subgroups-split-torus}
Every one-parameter subgroup
\[
\lambda:\Gm\longrightarrow\Gm^n
\]
has the form
\[
\lambda(t)
=
(t^{u_1},\ldots,t^{u_n})
\]
for a unique vector
\[
u=(u_1,\ldots,u_n)\in\ZZ^n.
\]
\end{theorem}

Indeed, composing \(\lambda\) with each coordinate projection
\[
\Gm^n\longrightarrow\Gm
\]
produces a character of \(\Gm\), and every such character is a power
\(t\mapsto t^{u_i}\).

\begin{figure}[ht]
\centering
\begin{tikzpicture}[
  mainbox/.style={
    rounded corners,
    draw=black!55,
    text width=3.1cm,
    minimum height=1.45cm,
    align=center,
    inner sep=5pt
  }
]
  \node[mainbox,fill=DictionaryBlue!8] (vector) at (-4.0,0)
    {\(\displaystyle u=(u_1,\ldots,u_n)\)\\[1mm]
     an integral vector};

  \node[mainbox,fill=PastelYellow] (formula) at (0,0)
    {\(\displaystyle
      (t^{u_1},\ldots,t^{u_n})\)\\[1mm]
     a monomial motion};

  \node[mainbox,fill=SurfaceGreen!9] (oneps) at (4.0,0)
    {\(\displaystyle \oneps{u}:\Gm\to\Gm^n\)\\[1mm]
     a one-parameter subgroup};

  \draw[very thick,-{Latex[length=2.4mm]},DictionaryBlue]
    (vector.east) -- (formula.west);
  \draw[very thick,-{Latex[length=2.4mm]},PredictionTeal]
    (formula.east) -- (oneps.west);
\end{tikzpicture}
\caption{After choosing coordinates, an integral vector and a
one-parameter subgroup are two descriptions of the same object.}
\label{fig:vector-monomial-one-parameter-subgroup}
\end{figure}

\subsection{The one-parameter subgroup lattice}

One-parameter subgroups can be multiplied pointwise:
\[
(\lambda\lambda')(t)
=
\lambda(t)\lambda'(t).
\]
For a coordinatized torus,
\[
\oneps{u}(t)\oneps{u'}(t)
=
\oneps{u+u'}(t).
\]
The inverse of \(\oneps{u}\) is \(\oneps{-u}\), and the identity is the
constant one-parameter subgroup
\[
\oneps{0}(t)=e.
\]

Thus \(X_*(T)\) is an abelian group.

\begin{theorem}
If \(T\) is an \(n\)-dimensional algebraic torus, then
\[
N:=X_*(T)
\]
is a lattice of rank \(n\).
\end{theorem}

For \(T=\Gm^n\), chosen coordinates give an identification
\[
N\simeq\ZZ^n.
\]
As with the character lattice, this identification depends on the
choice of coordinates, while the lattice \(N\) itself is intrinsic.

We will usually call \(N\) the \emph{one-parameter subgroup lattice}.
The notation is chosen to match the standard conventions of toric
geometry:
\[
M=X^*(T),
\qquad
N=X_*(T).
\]

\begin{observation}
The word \emph{cocharacter} is common in the theory of algebraic
groups. In toric geometry, the phrase \emph{one-parameter subgroup}
often carries more geometric intuition, because the map
\(\Gm\to T\) can be followed toward the boundary of a compactification.
\end{observation}

\subsection{One-parameter subgroups as motions}

A one-parameter subgroup always passes through the identity:
\[
\lambda(1)=e.
\]
To obtain a motion through an arbitrary point \(p\in T\), translate
the curve:
\[
t\longmapsto p\lambda(t).
\]
For \(T=\Gm^n\), if
\[
p=(a_1,\ldots,a_n)
\]
and
\[
\lambda=\oneps{u},
\]
then the translated motion is
\[
t\longmapsto
(a_1t^{u_1},\ldots,a_ns^{u_n}).
\]

The image of \(\lambda\) remains inside \(T\) as long as \(t\neq0\).
The interesting question is what happens as \(t\to0\). Since \(0\) is
not a point of \(\Gm\), this is really a question about whether the map
extends across \(t=0\) after \(T\) has been placed inside a larger
variety.

\begin{pauseandthink}
Why is the phrase ``take the limit as \(t\to0\)'' not intrinsic to the
torus \(T\) alone? What extra geometric object must be chosen before
such a limit can be discussed?
\end{pauseandthink}

A torus contains no boundary point at which the motion could arrive.
Limits become meaningful only after choosing an embedding
\[
T\subseteq X
\]
into a larger variety.

\subsection{Limits in affine space}

Consider the standard inclusion
\[
\Gm^n\subseteq\AAA^n.
\]
The one-parameter subgroup
\[
\oneps{u}(t)
=
(t^{u_1},\ldots,t^{u_n})
\]
extends to a morphism
\[
\AAA^1\longrightarrow\AAA^n
\]
if and only if
\[
u_i\geq0
\qquad
\text{for every }i.
\]
Indeed, a negative exponent produces a pole at \(t=0\).

When all \(u_i\geq0\), the limit is obtained coordinate by coordinate:
\[
\lim_{t\to0}t^{u_i}
=
\begin{cases}
1,&u_i=0,\\
0,&u_i>0.
\end{cases}
\]
Thus the vanishing pattern of the exponents determines the coordinate
stratum reached by the motion.

\begin{example}
Inside \(\Gm^2\subseteq\AAA^2\), consider:
\[
\begin{aligned}
\oneps{(2,1)}(t)&=(t^2,t),\\
\oneps{(0,3)}(t)&=(1,t^3),\\
\oneps{(-1,2)}(t)&=(t^{-1},t^2).
\end{aligned}
\]
The first motion extends to \(t=0\) and reaches the origin. The second
extends and reaches \((1,0)\). The third does not extend to
\(\AAA^1\to\AAA^2\), because its first coordinate has a pole.
\end{example}

For a translated motion
\[
t\longmapsto(a_1t^{u_1},\ldots,a_ns^{u_n}),
\]
the same criterion applies. When \(u_i=0\), the \(i\)-th coordinate
survives as \(a_i\); when \(u_i>0\), it tends to zero.

\begin{example}
Let \((a,b)\in\Gm^2\). The translated motion
\[
t\longmapsto(a,bs^3)
\]
approaches
\[
(a,0).
\]
As \(a\) varies, these motions reach every point of
\[
V(y)\setminus\{(0,0)\}.
\]
Thus one exponent vector describes a whole family of limits after
translation.
\end{example}

\begin{pauseandthink}
For \(u=(u_1,u_2)\), classify the limit of
\[
(a,b)\oneps{u}(t)
=
(as^{u_1},bs^{u_2})
\]
inside \(\AAA^2\), whenever the limit exists. Which choices of \(u\)
reach the open torus orbit, an open coordinate axis, or the origin?
\end{pauseandthink}

The answer reproduces the orbit stratification of \(\AAA^2\):

\[
\begin{array}{c|c}
\text{conditions on }u
&
\text{limit of }(a,b)\oneps{u}(t)
\\
\hline
u_1=u_2=0
&
(a,b)\in\Gm^2
\\
u_1>0,\ u_2=0
&
(0,b)\in V(x)\setminus\{0\}
\\
u_1=0,\ u_2>0
&
(a,0)\in V(y)\setminus\{0\}
\\
u_1>0,\ u_2>0
&
(0,0).
\end{array}
\]

\subsection{The same motion in a different compactification}

A one-parameter subgroup that has no limit in \(\AAA^n\) may acquire a
limit in another variety containing the same torus.

For the dense torus of \(\PP^2\), write
\[
T
=
\{[1:t_1:t_2]:t_1,t_2\in\Bbbk^\times\}.
\]
The one-parameter subgroup determined by \(u=(u_1,u_2)\) traces
\[
t\longmapsto[1:t^{u_1}:t^{u_2}].
\]
Even when one of the exponents is negative, projective rescaling may
produce a limit.

Set
\[
c=\min\{0,u_1,u_2\}.
\]
Then
\[
[1:t^{u_1}:t^{u_2}]
=
[t^{-c}:t^{u_1-c}:t^{u_2-c}],
\]
and all three exponents on the right are nonnegative. As \(t\to0\),
the coordinates whose original exponents attain the minimum survive.

\begin{example}
We obtain:
\[
\begin{aligned}
u=(1,2)
&\quad\Longrightarrow\quad
[1:t:t^2]\longrightarrow[1:0:0]=p_0,
\\
u=(-1,0)
&\quad\Longrightarrow\quad
[1:t^{-1}:1]=[t:1:t]\longrightarrow[0:1:0]=p_1,
\\
u=(0,-1)
&\quad\Longrightarrow\quad
[1:1:t^{-1}]=[t:t:1]\longrightarrow[0:0:1]=p_2,
\\
u=(-1,-1)
&\quad\Longrightarrow\quad
[1:t^{-1}:t^{-1}]=[t:1:1]\longrightarrow[0:1:1].
\end{aligned}
\]
The last limit lies in the one-dimensional orbit \(O_0\), rather than
at a fixed point.
\end{example}

This is a first indication that different regions of
\[
N_{\RR}=N\otimes_{\ZZ}\RR
\]
should correspond to different boundary limits. Cones in
\(N_{\RR}\) will eventually organize precisely this behavior.

\begin{prediction}
For the compactification \(T\subseteq\PP^2\), which vectors
\(u\in N_{\RR}\) lead to the same limiting orbit? Try to divide the
plane into regions according to the answer.
\end{prediction}

\subsection{Primitive one-parameter subgroups}

A nonzero lattice vector \(u\in N\) is called \emph{primitive} if it
is not a positive integer multiple of another lattice vector. Thus
\[
u=ku_0,\qquad k\geq1,
\]
with \(u_0\) primitive.

The one-parameter subgroups
\[
\oneps{u}
\qquad\text{and}\qquad
\oneps{u_0}
\]
have the same image subgroup, but they parametrize it differently:
\[
\oneps{u}(t)
=
\oneps{u_0}(t^k).
\]
The primitive vector is therefore the most economical lattice
generator of that one-dimensional direction.

This distinction becomes important when rays appear. A ray contains
many lattice points, but it has a unique primitive lattice generator.

\begin{figure}[ht]
\centering
\begin{tikzpicture}[scale=1.0]
  \draw[->,black!55] (-0.6,0) -- (6.2,0);
  \draw[->,black!55] (0,-0.6) -- (0,4.5);

  \draw[very thick,PredictionTeal,-{Latex[length=2.5mm]}]
    (0,0) -- (5.4,3.6);

  \foreach \x/\y/\lab in {
    1.2/0.8/{u_0},
    2.4/1.6/{2u_0},
    3.6/2.4/{3u_0},
    4.8/3.2/{4u_0}
  }{
    \fill[DictionaryBlue] (\x,\y) circle (2.2pt);
    \node[above left] at (\x,\y) {\(\lab\)};
  }

  \node[below right,PredictionTeal] at (4.8,3.2)
    {one ray, one primitive generator};
\end{tikzpicture}
\caption{All positive multiples of \(u_0\) lie on the same ray, but
\(u_0\) is its unique primitive lattice generator.}
\label{fig:primitive-generator-ray-preview}
\end{figure}

\subsection{Functoriality}

Let
\[
\phi:T_1\longrightarrow T_2
\]
be a morphism of algebraic tori. A one-parameter subgroup of \(T_1\)
can be followed by \(\phi\):
\[
\Gm\overset{\lambda}{\longrightarrow}T_1
\overset{\phi}{\longrightarrow}T_2.
\]
Hence \(\phi\) induces a homomorphism
\[
\phi_*:X_*(T_1)\longrightarrow X_*(T_2),
\qquad
\lambda\longmapsto\phi\circ\lambda.
\]

This time the direction is preserved:
\[
T_1\longrightarrow T_2
\qquad\text{produces}\qquad
X_*(T_1)\longrightarrow X_*(T_2).
\]

\begin{example}
Consider the multiplication map
\[
\mu:\Gm^2\longrightarrow\Gm,
\qquad
(t_1,t_2)\longmapsto t_1t_2.
\]
For \(u=(a,b)\in X_*(\Gm^2)\simeq\ZZ^2\),
\[
(\mu\circ\oneps{u})(t)
=
t^as^b
=
t^{a+b}.
\]
Thus
\[
\mu_*:\ZZ^2\longrightarrow\ZZ,
\qquad
(a,b)\longmapsto a+b.
\]
\end{example}

Characters and one-parameter subgroups therefore behave in opposite
ways under a torus morphism:
\[
\begin{array}{c@{\qquad}c}
\textbf{characters}
&
\textbf{one-parameter subgroups}
\\[1mm]
\phi^*:X^*(T_2)\to X^*(T_1)
&
\phi_*:X_*(T_1)\to X_*(T_2)
\\
\text{pull back}
&
\text{push forward}.
\end{array}
\]

\subsection{A first meeting between
\texorpdfstring{\(M\)}{M} and \texorpdfstring{\(N\)}{N}}

Take a character
\[
\chi:T\longrightarrow\Gm
\]
and a one-parameter subgroup
\[
\lambda:\Gm\longrightarrow T.
\]
Their composition is a character of \(\Gm\):
\[
\chi\circ\lambda:\Gm\longrightarrow\Gm.
\]
Therefore there is a unique integer \(r\) such that
\[
(\chi\circ\lambda)(t)=t^r.
\]

This integer has two equivalent interpretations:

\[
r
=
\text{the exponent of the character }\chi\circ\lambda
=
\ord_{t=0}(\chi\circ\lambda),
\]
where the second expression regards \(t^r\) as a rational function on
\(\AAA^1\). A positive value gives a zero, a negative value gives a
pole, and zero gives neither.

For coordinatized tori, if
\[
\chi=\character{m},
\qquad
m=(m_1,\ldots,m_n),
\]
and
\[
\lambda=\oneps{u},
\qquad
u=(u_1,\ldots,u_n),
\]
then
\[
(\character{m}\circ\oneps{u})(t)
=
t^{m_1u_1+\cdots+m_nu_n}.
\]

The next section will isolate this integer as the natural pairing
between the two lattices \(M\) and \(N\).

\begin{threedictionaries}
A one-parameter subgroup has three compatible interpretations.

\medskip

\noindent
\begin{minipage}[t]{0.30\linewidth}
\centering
\textbf{Algebra}\par
\smallskip
It is a homomorphism
\[
\lambda:\Gm\to T.
\]
Such maps form the lattice
\[
N=X_*(T).
\]
\end{minipage}
\hfill
\begin{minipage}[t]{0.30\linewidth}
\centering
\textbf{Geometry}\par
\smallskip
It is an algebraic motion inside \(T\). In a compactification, its
closure may meet a boundary orbit.
\end{minipage}
\hfill
\begin{minipage}[t]{0.30\linewidth}
\centering
\textbf{Combinatorics}\par
\smallskip
After choosing coordinates, it is an integral vector
\[
u\in\ZZ^n.
\]
Rays remember primitive directions.
\end{minipage}
\end{threedictionaries}

\begin{guidingquestion}
Why do characters and one-parameter subgroups form dual lattices, and
what geometric information is measured by their pairing?
\end{guidingquestion}


\section{Dual lattices and the pairing}
\label{sec:dual-lattices-pairing}
\index{dual lattice}
\index{lattice pairing}
\index{order of vanishing}

We now have two lattices attached to a torus:
\[
M=X^*(T)
=
\Hom_{\mathrm{alg.grp}}(T,\Gm)
\]
and
\[
N=X_*(T)
=
\Hom_{\mathrm{alg.grp}}(\Gm,T).
\]
The arrows point in opposite directions, so the two kinds of maps can
be composed:
\[
\Gm
\overset{\lambda}{\longrightarrow}
T
\overset{\chi}{\longrightarrow}
\Gm.
\]
This composition is again a character of \(\Gm\), and therefore it is
an integral power map. The exponent of that power is the basic
numerical invariant connecting \(M\) and \(N\).

\subsection{The pairing by composition}

Let
\[
m\in M
\qquad\text{and}\qquad
u\in N.
\]
We write the corresponding character and one-parameter subgroup as
\[
\character{m}:T\longrightarrow\Gm,
\qquad
\oneps{u}:\Gm\longrightarrow T.
\]
Their composition is a character of \(\Gm\), so there is a unique
integer \(\pair{m}{u}\) such that
\[
\character{m}\circ\oneps{u}(t)
=
t^{\pair{m}{u}}.
\]

\begin{definition}
The integer
\[
\pair{m}{u}\in\ZZ
\]
is the \emph{pairing} of the character \(m\) with the one-parameter
subgroup \(u\).
\end{definition}

The pairing may be read in several equivalent ways:
\[
\pair{m}{u}
=
\text{the exponent of }\character{m}\circ\oneps{u}
=
\ord_{t=0}\bigl(\character{m}\circ\oneps{u}\bigr).
\]
Thus

\[
\begin{array}{c|c}
\pair{m}{u}>0
&
\character{m}\circ\oneps{u}
\text{ has a zero of order }\pair{m}{u}\text{ at }0,
\\
\pair{m}{u}=0
&
\character{m}\circ\oneps{u}
\text{ has neither a zero nor a pole at }0,
\\
\pair{m}{u}<0
&
\character{m}\circ\oneps{u}
\text{ has a pole of order }-\pair{m}{u}\text{ at }0.
\end{array}
\]

This is the intrinsic version of the coordinate calculation previewed
in the previous section.

\begin{example}
Let \(T=\Gm^3\), and take
\[
m=(2,-1,4)\in M\simeq\ZZ^3,
\qquad
u=(-1,3,2)\in N\simeq\ZZ^3.
\]
Then
\[
\character{m}(t_1,t_2,t_3)
=
t_1^2t_2^{-1}t_3^4
\]
and
\[
\oneps{u}(t)
=
(t^{-1},t^3,t^2).
\]
Therefore
\[
\begin{aligned}
\character{m}\bigl(\oneps{u}(t)\bigr)
&=
(t^{-1})^2(t^3)^{-1}(t^2)^4\\
&=
t^{-2-3+8}\\
&=
t^3.
\end{aligned}
\]
Hence
\[
\pair{m}{u}=3.
\]
\end{example}

\begin{pauseandthink}
Compute \(\pair{m}{u}\) in the previous example without composing the
maps explicitly. What familiar operation on the two coordinate
vectors appears?
\end{pauseandthink}

After choosing dual coordinate bases, the pairing becomes the usual
dot product:
\[
\boxed{
\pair{m}{u}
=
m_1u_1+\cdots+m_nu_n.
}
\]
The formula is convenient, but the pairing itself was defined without
coordinates.

\subsection{Why the pairing is bilinear}

The pairing is additive in each variable:
\[
\pair{m+m'}{u}
=
\pair{m}{u}+\pair{m'}{u}
\]
and
\[
\pair{m}{u+u'}
=
\pair{m}{u}+\pair{m}{u'}.
\]

The first identity follows because multiplication of characters
corresponds to addition in \(M\):
\[
\begin{aligned}
\character{m+m'}\circ\oneps{u}
&=
(\character{m}\character{m'})\circ\oneps{u}\\
&=
(\character{m}\circ\oneps{u})
(\character{m'}\circ\oneps{u})\\
&=
t^{\pair{m}{u}}
t^{\pair{m'}{u}}.
\end{aligned}
\]
The exponents add. The second identity follows similarly from
\[
\oneps{u+u'}(t)
=
\oneps{u}(t)\oneps{u'}(t).
\]

Thus we have a bilinear map of lattices
\[
\pair{\phantom{m}}{\phantom{u}}:
M\times N\longrightarrow\ZZ.
\]

\begin{observation}
The word \emph{bilinear} here means additive in each variable because
\(M\) and \(N\) are abelian groups. After tensoring with \(\RR\), the
same pairing becomes bilinear in the usual real-vector-space sense.
\end{observation}

\subsection{Perfect duality}

Fix \(m\in M\). The rule
\[
u\longmapsto\pair{m}{u}
\]
is a homomorphism
\[
N\longrightarrow\ZZ.
\]
Therefore every character determines an element of the dual lattice
\[
N^\vee
=
\Hom_{\ZZ}(N,\ZZ).
\]
This gives a homomorphism
\[
M\longrightarrow N^\vee.
\]

Likewise, fixing \(u\in N\) gives a homomorphism
\[
M\longrightarrow\ZZ,
\qquad
m\longmapsto\pair{m}{u},
\]
and therefore a map
\[
N\longrightarrow M^\vee.
\]

\begin{theorem}[Duality of characters and one-parameter subgroups]
\label{thm:character-oneps-duality}
For an algebraic torus \(T\), the pairing is perfect. Equivalently,
the natural maps
\[
M\overset{\sim}{\longrightarrow}
\Hom_{\ZZ}(N,\ZZ)
\]
and
\[
N\overset{\sim}{\longrightarrow}
\Hom_{\ZZ}(M,\ZZ)
\]
are isomorphisms.
\end{theorem}

Thus \(M\) and \(N\) are dual lattices:
\[
\boxed{
M=N^\vee,
\qquad
N=M^\vee.
}
\]

For \(T=\Gm^n\), this statement is visible in coordinates. Let
\[
e_1,\ldots,e_n
\]
be the coordinate character basis of \(M\), and let
\[
e_1^\vee,\ldots,e_n^\vee
\]
be the coordinate one-parameter subgroup basis of \(N\). Then
\[
\pair{e_i}{e_j^\vee}
=
\delta_{ij}.
\]
These are dual bases.

\begin{figure}[ht]
\centering
\begin{tikzpicture}[
  mainbox/.style={
    rounded corners,
    draw=black!55,
    text width=3.35cm,
    minimum height=1.55cm,
    align=center,
    inner sep=5pt
  }
]
  \node[mainbox,fill=DictionaryBlue!8] (M) at (-4.0,0)
    {\(\displaystyle M=X^*(T)\)\\[1mm]
     characters\\
     \(T\to\Gm\)};

  \node[mainbox,fill=PastelYellow] (Z) at (0,0)
    {\(\displaystyle \pair{m}{u}\in\ZZ\)\\[1mm]
     exponent of\\
     \(\character{m}\circ\oneps{u}\)};

  \node[mainbox,fill=SurfaceGreen!9] (N) at (4.0,0)
    {\(\displaystyle N=X_*(T)\)\\[1mm]
     one-parameter subgroups\\
     \(\Gm\to T\)};

  \draw[very thick,-{Latex[length=2.4mm]},DictionaryBlue]
    (M.east) -- (Z.west);
  \draw[very thick,-{Latex[length=2.4mm]},PredictionTeal]
    (N.west) -- (Z.east);
\end{tikzpicture}
\caption{A character and a one-parameter subgroup compose to an
integral power map on \(\Gm\). This composition produces the perfect
pairing between \(M\) and \(N\).}
\label{fig:character-oneps-pairing}
\end{figure}

\begin{pauseandthink}
Suppose \(m\in M\) satisfies
\[
\pair{m}{u}=0
\qquad
\text{for every }u\in N.
\]
Why should \(m=0\)? Interpret this first in coordinates, and then in
terms of the perfect-pairing theorem.
\end{pauseandthink}

\subsection{Real vector spaces and linear functions}

Tensoring with \(\RR\) gives dual real vector spaces
\[
M_{\RR}=M\otimes_{\ZZ}\RR,
\qquad
N_{\RR}=N\otimes_{\ZZ}\RR.
\]
The integral pairing extends uniquely to a real bilinear pairing
\[
\pair{\phantom{m}}{\phantom{u}}:
M_{\RR}\times N_{\RR}\longrightarrow\RR.
\]
Consequently,
\[
M_{\RR}\simeq N_{\RR}^{\vee},
\qquad
N_{\RR}\simeq M_{\RR}^{\vee}.
\]

A point \(m\in M_{\RR}\) may therefore be viewed as a linear function
on \(N_{\RR}\):
\[
u\longmapsto\pair{m}{u}.
\]
Similarly, a point \(u\in N_{\RR}\) defines a linear function on
\(M_{\RR}\).

This is the reason inequalities in toric geometry naturally take the
form
\[
\pair{m}{u}\geq c.
\]
The vector \(u\) determines a linear functional on \(M_{\RR}\), and
the inequality cuts out a half-space.

\begin{figure}[ht]
\centering
\begin{tikzpicture}[scale=1.0]
  \draw[->,black!55] (-3.4,0) -- (3.6,0)
    node[right] {\(N_{\RR}\)};
  \draw[->,black!55] (0,-2.8) -- (0,3.0);

  \draw[very thick,PredictionTeal]
    (-2.8,-1.4) -- (2.8,1.4);
  \node[PredictionTeal,above left] at (2.5,1.25)
    {\(\pair{m}{u}=0\)};

  \fill[DictionaryBlue] (1.1,2.2) circle (2.2pt);
  \node[above right] at (1.1,2.2) {\(u_+\)};
  \node[align=left] at (2.15,2.25)
    {\(\pair{m}{u_+}>0\)};

  \fill[WarningRed] (-1.1,-2.2) circle (2.2pt);
  \node[below left] at (-1.1,-2.2) {\(u_-\)};
  \node[align=right] at (-2.15,-2.25)
    {\(\pair{m}{u_-}<0\)};

  \node[align=center] at (-2.0,1.45)
    {positive\\half-space};
  \node[align=center] at (2.0,-1.45)
    {negative\\half-space};
\end{tikzpicture}
\caption{A character \(m\in M_{\RR}\) acts as a linear function on
\(N_{\RR}\). Its zero set is a hyperplane separating positive and
negative regions.}
\label{fig:character-linear-functional-on-N}
\end{figure}

\begin{observation}
In dimension two, the zero set of a nonzero linear functional is a
line through the origin. In higher dimensions it is a hyperplane.
This linear-algebraic picture will organize cones, dual cones, and
polyhedra throughout the course.
\end{observation}

\subsection{Orders along one-parameter motions}

The pairing has a direct geometric meaning. Suppose \(T\) is embedded
as a dense open subset of a variety \(X\), and suppose the
one-parameter subgroup \(\oneps{u}\) extends to a map
\[
\overline{\oneps{u}}:\AAA^1\longrightarrow X.
\]
Then a character \(\character{m}\), viewed as a rational function on
\(X\), pulls back to
\[
\overline{\oneps{u}}^{\,*}(\character{m})
=
t^{\pair{m}{u}}.
\]
Hence
\[
\pair{m}{u}
=
\ord_{t=0}
\bigl(\overline{\oneps{u}}^{\,*}(\character{m})\bigr).
\]

This is different from the coordinate-hyperplane interpretation of an
individual exponent \(m_i\). There, after choosing a specific
compactification
\[
\Gm^n\subseteq\AAA^n,
\]
the integer \(m_i\) measures the order of \(\character{m}\) along the
coordinate divisor \(V(x_i)\). Here, for an arbitrary one-parameter
subgroup \(u\), the integer \(\pair{m}{u}\) measures the order seen
along the motion determined by \(u\).

The two viewpoints agree in the coordinate case. If
\[
u=e_i^\vee,
\]
then
\[
\pair{m}{e_i^\vee}=m_i.
\]

\begin{example}
In \(T=\Gm^2\), let
\[
\character{m}(x,y)=x^2y^{-3}
\]
and let
\[
\oneps{u}(t)=(t^4,t).
\]
Then
\[
\pair{m}{u}
=
2\cdot4-3\cdot1
=
5.
\]
Along this motion,
\[
\character{m}(t^4,t)=t^5,
\]
so the character vanishes to order \(5\) at the limiting parameter
\(t=0\).
\end{example}

\begin{prediction}
Suppose a collection of characters
\[
\character{m_1},\ldots,\character{m_r}
\]
is evaluated along the same one-parameter subgroup \(u\). How might
the numbers
\[
\pair{m_1}{u},\ldots,\pair{m_r}{u}
\]
help predict which monomials dominate, vanish, or survive as
\(t\to0\)?
\end{prediction}

\subsection{Compatibility with torus morphisms}

Let
\[
\phi:T_1\longrightarrow T_2
\]
be a morphism of algebraic tori. Recall that it induces
\[
\phi^*:M_2\longrightarrow M_1
\]
on character lattices and
\[
\phi_*:N_1\longrightarrow N_2
\]
on one-parameter subgroup lattices.

These two maps are adjoint with respect to the pairings:
\[
\boxed{
\pair{\phi^*(m)}{u}
=
\pair{m}{\phi_*(u)}.
}
\]
Indeed, both sides are the exponent of the same composition:
\[
\Gm
\overset{\oneps{u}}{\longrightarrow}
T_1
\overset{\phi}{\longrightarrow}
T_2
\overset{\character{m}}{\longrightarrow}
\Gm.
\]

\begin{example}
Let
\[
\phi:\Gm^2\longrightarrow\Gm^2,
\qquad
(t_1,t_2)\longmapsto(t_1^2t_2,t_1t_2^3).
\]
On character lattices,
\[
\phi^*(a,b)
=
(2a+b,a+3b).
\]
On one-parameter subgroup lattices,
\[
\phi_*(p,q)
=
(2p+q,p+3q).
\]
For
\[
m=(a,b),
\qquad
u=(p,q),
\]
one checks that
\[
\pair{\phi^*(m)}{u}
=
(2a+b)p+(a+3b)q
\]
and
\[
\pair{m}{\phi_*(u)}
=
a(2p+q)+b(p+3q),
\]
which are equal.
\end{example}

In this example, the numerical matrices representing \(\phi^*\) and
\(\phi_*\) happen to agree because the matrix is symmetric. In
general, the two matrices are transposes of one another once dual
bases are fixed. We will return to this carefully in the next
section.

\begin{threedictionaries}
The pairing is the bridge between the two lattices.

\medskip

\noindent
\begin{minipage}[t]{0.30\linewidth}
\centering
\textbf{Algebra}\par
\smallskip
Composition gives
\[
\character{m}\circ\oneps{u}
=
t^{\pair{m}{u}}.
\]
The pairing is bilinear and perfect.
\end{minipage}
\hfill
\begin{minipage}[t]{0.30\linewidth}
\centering
\textbf{Geometry}\par
\smallskip
\(\pair{m}{u}\) is the order of vanishing or pole of a character along
a one-parameter motion.
\end{minipage}
\hfill
\begin{minipage}[t]{0.30\linewidth}
\centering
\textbf{Combinatorics}\par
\smallskip
\[
M=N^\vee,
\qquad
N=M^\vee.
\]
Elements of one lattice are integral linear functions on the other.
\end{minipage}
\end{threedictionaries}

\begin{guidingquestion}
How do the lattices, the pairing, and the formulas for torus morphisms
change when we replace one basis by another?
\end{guidingquestion}


\section{Changing coordinates on a torus}
\label{sec:changing-torus-coordinates}
\index{change of lattice basis}
\index{monomial automorphism}
\index{general linear group over the integers}

A torus does not come with preferred coordinates. Choosing coordinates
means choosing a basis of its character lattice, and changing
coordinates means replacing that basis by another one. The allowable
changes are therefore controlled by
\[
\operatorname{GL}(n,\ZZ),
\]
the group of invertible \(n\times n\) integer matrices.

This is an important difference between affine space and a torus.
Coordinate changes on \(\AAA^n\) may involve many kinds of
polynomials. Coordinate changes that preserve the algebraic group
structure of \(\Gm^n\) are monomial.

\subsection{From a new basis to new torus coordinates}

Let \(T\) be an \(n\)-dimensional torus, and choose a basis
\[
m_1,\ldots,m_n
\]
of its character lattice \(M\). The corresponding torus coordinates
are
\[
x_i=\character{m_i},
\qquad
i=1,\ldots,n.
\]

Now choose another basis
\[
m'_1,\ldots,m'_n.
\]
Write the new basis vectors in terms of the old basis:
\[
m'_j=\sum_{i=1}^n a_{ij}m_i.
\]
The coefficients form an integer matrix
\[
A=(a_{ij}).
\]
Because both collections are bases of the same lattice,
\[
A\in\operatorname{GL}(n,\ZZ).
\]

The new coordinate characters are
\[
x'_j
=
\character{m'_j}
=
\prod_{i=1}^n x_i^{a_{ij}}.
\]
Thus every change of basis gives an invertible monomial change of
coordinates:
\[
(x_1,\ldots,x_n)
\longmapsto
(x'_1,\ldots,x'_n).
\]

\begin{example}
Let \(T=\Gm^2\), and begin with the standard character basis
\[
m_1=(1,0),
\qquad
m_2=(0,1).
\]
Choose the new basis
\[
m'_1=-m_1,
\qquad
m'_2=-m_1+m_2.
\]
Its change-of-basis matrix is
\[
A=
\begin{pmatrix}
-1&-1\\
0&1
\end{pmatrix},
\]
whose columns record \(m'_1\) and \(m'_2\) in the old basis.

If
\[
x=\character{m_1},
\qquad
y=\character{m_2},
\]
then
\[
x'=x^{-1},
\qquad
y'=x^{-1}y.
\]
These are exactly the coordinate changes that appeared between the
charts \(U_0\) and \(U_1\) on the dense torus of \(\PP^2\).
\end{example}

\begin{pauseandthink}
Why must the matrix \(A\) have determinant \(\pm1\), rather than merely
a nonzero determinant? What would go wrong if the columns generated a
proper finite-index sublattice of \(M\)?
\end{pauseandthink}

If \(|\det A|>1\), the new vectors do not form a lattice basis. The
associated monomial map may still be a finite morphism of tori, but it
is not an isomorphism.

\subsection{Monomial automorphisms of a split torus}

Let
\[
A=(a_{ij})\in\operatorname{GL}(n,\ZZ).
\]
It defines a map
\[
\Phi_A:\Gm^n\longrightarrow\Gm^n
\]
by
\[
\Phi_A(t_1,\ldots,t_n)
=
\left(
\prod_{i=1}^n t_i^{a_{i1}},
\ldots,
\prod_{i=1}^n t_i^{a_{in}}
\right).
\]
The columns of \(A\) contain the exponent vectors of the output
coordinates.

Because \(A^{-1}\) also has integer entries, the map
\[
\Phi_{A^{-1}}
\]
is again monomial and is inverse to \(\Phi_A\). Hence \(\Phi_A\) is an
automorphism of algebraic tori.

Conversely, every algebraic group automorphism of \(\Gm^n\) arises in
this way. Therefore
\[
\Aut_{\mathrm{alg.grp}}(\Gm^n)
\simeq
\operatorname{GL}(n,\ZZ).
\]

\begin{observation}
Translations
\[
t\longmapsto p\,t,
\qquad
p\in T,
\]
are automorphisms of the underlying variety, but they do not preserve
the identity unless \(p=e\). They are therefore not automorphisms of
the algebraic group. The group automorphisms are precisely the
monomial coordinate changes coming from lattice automorphisms.
\end{observation}

\begin{figure}[ht]
\centering
\begin{tikzpicture}[
  mainbox/.style={
    rounded corners,
    draw=black!55,
    text width=3.35cm,
    minimum height=1.55cm,
    align=center,
    inner sep=5pt
  }
]
  \node[mainbox,fill=DictionaryBlue!8] (matrix) at (-4.0,0)
    {\(\displaystyle A\in\operatorname{GL}(n,\ZZ)\)\\[1mm]
     a lattice automorphism};

  \node[mainbox,fill=PastelYellow] (basis) at (0,0)
    {a change of basis\\[1mm]
     \(\displaystyle m'_j=\sum_i a_{ij}m_i\)};

  \node[mainbox,fill=SurfaceGreen!9] (map) at (4.0,0)
    {an invertible monomial map\\[1mm]
     \(\displaystyle
     x'_j=\prod_i x_i^{a_{ij}}\)};

  \draw[very thick,<->,DictionaryBlue]
    (matrix.east) -- (basis.west);
  \draw[very thick,<->,PredictionTeal]
    (basis.east) -- (map.west);
\end{tikzpicture}
\caption{Integer change-of-basis matrices, lattice automorphisms, and
monomial automorphisms of a torus are three descriptions of the same
operation.}
\label{fig:glz-basis-monomial-automorphism}
\end{figure}

\subsection{What happens to the dual basis?}

Let
\[
m_1,\ldots,m_n
\]
be a basis of \(M\), and let
\[
u_1,\ldots,u_n
\]
be the dual basis of \(N\):
\[
\pair{m_i}{u_j}=\delta_{ij}.
\]
Suppose the new character basis is given by
\[
m'_j=\sum_i a_{ij}m_i.
\]
Let
\[
u'_1,\ldots,u'_n
\]
be the basis of \(N\) dual to the \(m'_j\).

In matrix notation, the new character basis is obtained from the old
one by \(A\). The new one-parameter subgroup basis is obtained by
\[
A^{-T}=(A^{-1})^T.
\]
Indeed, if the columns of a matrix \(B\) record the \(u'_j\) in the old
basis, the duality condition says
\[
A^TB=I,
\]
and therefore
\[
B=A^{-T}.
\]

Thus:

\[
\boxed{
\text{characters change by }A,
\qquad
\text{one-parameter subgroups change by }A^{-T}.
}
\]

\begin{example}
For
\[
A=
\begin{pmatrix}
-1&-1\\
0&1
\end{pmatrix},
\]
one has
\[
A^{-1}=A
\qquad\text{and}\qquad
A^{-T}=
\begin{pmatrix}
-1&0\\
-1&1
\end{pmatrix}.
\]
Hence the basis of \(N\) dual to
\[
m'_1=-m_1,
\qquad
m'_2=-m_1+m_2
\]
is
\[
u'_1=-u_1-u_2,
\qquad
u'_2=u_2.
\]
A direct computation verifies
\[
\pair{m'_i}{u'_j}=\delta_{ij}.
\]
\end{example}

\begin{pauseandthink}
Why does the inverse transpose appear? Start with the requirement
\[
\pair{m'_i}{u'_j}=\delta_{ij}
\]
and translate it into a matrix equation.
\end{pauseandthink}

\subsection{Coordinates of vectors versus bases}

There are two related matrix operations that are easy to confuse.

First, the basis vectors themselves change:
\[
(m'_1,\ldots,m'_n)
=
(m_1,\ldots,m_n)A.
\]
Second, the coordinate column of a fixed element changes in the
opposite direction.

If
\[
m=\sum_i c_i m_i=\sum_j c'_j m'_j,
\]
then
\[
c=Ac',
\qquad
c'=A^{-1}c.
\]
Thus the basis changes by \(A\), while the coordinate column of a fixed
character changes by \(A^{-1}\).

Similarly, if
\[
u=\sum_i d_i u_i=\sum_j d'_j u'_j,
\]
then
\[
d=A^{-T}d',
\qquad
d'=A^Td.
\]

The pairing is unchanged:
\[
c^Td
=
(c')^Td'.
\]
Indeed,
\[
(Ac')^T(A^{-T}d')
=
(c')^Td'.
\]

\begin{warningbox}
A matrix may describe either a change of basis or the coordinates of a
linear map. Transposes and inverses depend on which convention is
being used. In these notes, columns record images of basis vectors.
Whenever a matrix formula matters, we will state the convention
explicitly.
\end{warningbox}

\subsection{Matrices for morphisms of tori}

Now let
\[
\phi:T_1\longrightarrow T_2
\]
be a morphism of tori. Choose bases
\[
m_1^{(1)},\ldots,m_r^{(1)}
\]
of \(M_1=X^*(T_1)\) and
\[
m_1^{(2)},\ldots,m_s^{(2)}
\]
of \(M_2=X^*(T_2)\).

The pullback
\[
\phi^*:M_2\longrightarrow M_1
\]
is represented by an integer \(r\times s\) matrix
\[
B=(b_{ij}),
\]
defined by
\[
\phi^*(m_j^{(2)})
=
\sum_{i=1}^r b_{ij}m_i^{(1)}.
\]
Equivalently, if
\[
x_i=\character{m_i^{(1)}}
\qquad\text{and}\qquad
y_j=\character{m_j^{(2)}},
\]
then
\[
\phi^*(y_j)
=
\prod_{i=1}^r x_i^{b_{ij}}.
\]
On points, the torus morphism is therefore
\[
\phi(t_1,\ldots,t_r)
=
\left(
\prod_i t_i^{b_{i1}},
\ldots,
\prod_i t_i^{b_{is}}
\right).
\]

The induced map on one-parameter subgroup lattices,
\[
\phi_*:N_1\longrightarrow N_2,
\]
is represented in the dual bases by the transpose matrix
\[
B^T.
\]
This follows from the adjointness identity
\[
\pair{\phi^*(m)}{u}
=
\pair{m}{\phi_*(u)}.
\]

\begin{example}
Consider
\[
\phi:\Gm^2\longrightarrow\Gm^2,
\qquad
(t_1,t_2)
\longmapsto
(t_1^2t_2,\ t_1t_2^3).
\]
The pullback of the two target coordinate characters is
\[
\phi^*(y_1)=x_1^2x_2,
\qquad
\phi^*(y_2)=x_1x_2^3.
\]
Thus
\[
B=
\begin{pmatrix}
2&1\\
1&3
\end{pmatrix}.
\]
In this example \(B=B^T\), so the same numerical matrix represents
\(\phi^*\) and \(\phi_*\). This is accidental: the matrix happens to
be symmetric.
\end{example}

\begin{example}
Consider instead
\[
\psi:\Gm^2\longrightarrow\Gm^2,
\qquad
(t_1,t_2)
\longmapsto
(t_1^2t_2^3,\ t_1^{-1}t_2^4).
\]
Then
\[
B=
\begin{pmatrix}
2&-1\\
3&4
\end{pmatrix}
\]
represents
\[
\psi^*:M_2\longrightarrow M_1,
\]
while
\[
B^T=
\begin{pmatrix}
2&3\\
-1&4
\end{pmatrix}
\]
represents
\[
\psi_*:N_1\longrightarrow N_2.
\]
For \(u=(p,q)\in N_1\),
\[
\psi\bigl(\oneps{u}(t)\bigr)
=
\bigl(t^{2p+3q},t^{-p+4q}\bigr),
\]
which indeed has exponent vector
\[
B^T
\begin{pmatrix}
p\\q
\end{pmatrix}.
\]
\end{example}

\subsection{Isomorphisms, finite maps, and dimensions}

The matrix description makes several basic properties visible.

If \(B\) is square and
\[
\det B=\pm1,
\]
then \(B\in\operatorname{GL}(n,\ZZ)\), and the corresponding torus
morphism is an isomorphism.

If \(B\) is square and
\[
\det B\neq0
\]
but \(|\det B|>1\), then the map is finite and surjective but not an
isomorphism. For example,
\[
\Gm\longrightarrow\Gm,
\qquad
t\longmapsto t^d
\]
corresponds to multiplication by \(d\) on the character lattice.

Over an algebraically closed field whose characteristic does not
divide \(d\), the kernel consists of the \(d\)-th roots of unity. In
positive characteristic dividing \(d\), the scheme-theoretic kernel
requires more care; we will not need that refinement here.

When the source and target have different dimensions, the ranks of
the lattice maps help describe kernels and images. We will use these
ideas later, but the main lesson for now is simpler:

\[
\boxed{
\text{morphisms of tori are encoded by homomorphisms of lattices.}
}
\]

\begin{guidingquestion}
Suppose a torus acts linearly on a finite-dimensional vector space.
After choosing coordinates on the torus, its elements are tuples of
nonzero scalars. What kind of basis of the vector space might make the
action as simple as possible?
\end{guidingquestion}

\begin{threedictionaries}
Changing coordinates on a torus is an integral operation.

\medskip

\noindent
\begin{minipage}[t]{0.30\linewidth}
\centering
\textbf{Algebra}\par
\smallskip
Torus morphisms pull back characters:
\[
\phi^*:M_2\to M_1.
\]
After choosing bases, this is an integer matrix.
\end{minipage}
\hfill
\begin{minipage}[t]{0.30\linewidth}
\centering
\textbf{Geometry}\par
\smallskip
Automorphisms of \(\Gm^n\) preserving the identity are invertible
monomial coordinate changes.
\end{minipage}
\hfill
\begin{minipage}[t]{0.30\linewidth}
\centering
\textbf{Combinatorics}\par
\smallskip
Changes of lattice basis lie in
\[
\operatorname{GL}(n,\ZZ).
\]
The dual lattice changes by the inverse transpose.
\end{minipage}
\end{threedictionaries}



\section{Representation theory in one page}
\label{sec:torus-representation-theory}
\index{representation!of a torus}
\index{weight space}
\index{weight-space decomposition}
\index{invariant vector}

Characters are the one-dimensional representations of a torus. The
basic representation theory of tori says that every finite-dimensional
representation is assembled from these one-dimensional pieces.

This is the central structural fact of the chapter. We will state it
without proof and focus on learning how to use it.

\subsection{What is a representation?}

Let \(V\) be a finite-dimensional vector space over \(\Bbbk\).

\begin{definition}
A \emph{finite-dimensional algebraic representation} of a torus \(T\)
on \(V\) is a morphism of algebraic groups
\[
\rho:T\longrightarrow\operatorname{GL}(V).
\]
Equivalently, it is an algebraic action
\[
T\times V\longrightarrow V
\]
such that, for each \(t\in T\), the map
\[
v\longmapsto t\cdot v
\]
is linear.
\end{definition}

The adjective \emph{algebraic} matters. It requires the matrix entries of \(\rho(t)\) to be regular functions of \(t\), after bases are chosen.

\begin{example}
Let \(T=\Gm\) act on \(V=\Bbbk^2\) by
\[
t\cdot(x,y)=(t^2x,t^{-3}y).
\]
In the standard basis,
\[
\rho(t)
=
\begin{pmatrix}
t^2&0\\
0&t^{-3}
\end{pmatrix}.
\]
The two coordinate lines are one-dimensional representations with
characters \(t\mapsto t^2\) and \(t\mapsto t^{-3}\).
\end{example}

\begin{example}
Let \(T=\Gm^2\) act on \(\Bbbk^3\) by
\[
(t_1,t_2)\cdot(x_0,x_1,x_2)
=
(x_0,t_1x_1,t_2x_2).
\]
The three coordinate lines transform by the characters
\[
1,\qquad t_1,\qquad t_2.
\]
Projectivizing this representation gives the torus action on
\(\PP^2\) used throughout Chapter~0:
\[
(t_1,t_2)\cdot[x_0:x_1:x_2]
=
[x_0:t_1x_1:t_2x_2].
\]
\end{example}

\subsection{Weight vectors and weight spaces}

A character tells us how a torus may scale a single vector.

\begin{definition}
Let \(V\) be a representation of \(T\), and let \(m\in M=X^*(T)\).
A nonzero vector \(v\in V\) is a \emph{weight vector of weight \(m\)}
if
\[
t\cdot v=\character{m}(t)v
\qquad
\text{for every }t\in T.
\]
The corresponding \emph{weight space} is
\[
V_m
=
\{v\in V:t\cdot v=\character{m}(t)v
\text{ for every }t\in T\}.
\]
\end{definition}

The zero vector lies in every \(V_m\), but a nonzero vector cannot
have two distinct weights.

\begin{example}
For the action
\[
t\cdot(x,y)=(t^2x,t^{-3}y)
\]
on \(\Bbbk^2\), let \(e_1,e_2\) be the standard basis. Then
\[
e_1\in V_2,
\qquad
e_2\in V_{-3},
\]
and
\[
V=V_2\oplus V_{-3}.
\]
The vector \(e_1+e_2\) is not a weight vector: its two components scale
by different characters.
\end{example}

\begin{pauseandthink}
For the action of \(\Gm\) on \(\Bbbk^3\) given by
\[
t\cdot(x,y,z)=(t^2x,t^2y,t^{-1}z),
\]
find all weight spaces. Which vectors are weight vectors, and which
are not?
\end{pauseandthink}

Here
\[
V_2=\Span\{e_1,e_2\},
\qquad
V_{-1}=\Span\{e_3\}.
\]
Every nonzero vector in either of these subspaces is a weight vector.
A vector having nonzero components in both subspaces is not.

This example is worth remembering: a weight space need not be
one-dimensional. A weight records how the torus acts, not the
dimension of the subspace on which it acts that way.

\subsection{The weight-space decomposition}

The representation theory of a torus is governed by the following
theorem.

\begin{theorem}[Weight-space decomposition]
\label{thm:torus-weight-decomposition}
Let \(T\) be an algebraic torus, and let \(V\) be a finite-dimensional
algebraic representation of \(T\). Then
\[
V
=
\bigoplus_{m\in M}V_m.
\]
Only finitely many of the weight spaces \(V_m\) are nonzero.
\end{theorem}

In other words, every vector can be written uniquely as a finite sum
of weight vectors belonging to distinct weights.

The finite set
\[
\operatorname{Wt}(V)
=
\{m\in M:V_m\neq0\}
\]
is called the \emph{set of weights} of the representation. The
dimension
\[
\dim V_m
\]
is the \emph{multiplicity} of the weight \(m\).

\begin{observation}
After choosing a basis of each nonzero weight space, every matrix
\(\rho(t)\) becomes diagonal:
\[
\rho(t)
=
\diag
\bigl(
\character{m_1}(t),\ldots,\character{m_r}(t)
\bigr),
\]
where weights may be repeated according to their multiplicities.

Thus a torus representation may be simultaneously diagonalized:
one basis diagonalizes the action of every element of \(T\) at once.
\end{observation}

\begin{figure}[ht]
\centering
\begin{tikzpicture}[
  weightbox/.style={
    rounded corners,
    draw=black!55,
    minimum width=2.55cm,
    minimum height=1.25cm,
    align=center,
    inner sep=4pt
  }
]
  \node at (0,1.75) {\(\displaystyle V=\bigoplus_{m\in\operatorname{Wt}(V)}V_m\)};

  \node[weightbox,fill=DictionaryBlue!8] (v1) at (-4.2,0)
    {\(V_{m_1}\)\\
     \(t\cdot v=\character{m_1}(t)v\)};

  \node[weightbox,fill=PastelYellow] (v2) at (-1.4,0)
    {\(V_{m_2}\)\\
     \(t\cdot v=\character{m_2}(t)v\)};

  \node at (1.0,0) {\(\cdots\)};

  \node[weightbox,fill=SurfaceGreen!9] (vr) at (3.5,0)
    {\(V_{m_r}\)\\
     \(t\cdot v=\character{m_r}(t)v\)};

  \draw[thick,black!45]
    (-5.25,-0.92) -- (4.55,-0.92);
  \node[below=5pt,align=center] at (-0.35,-0.92)
    {each summand is scaled by one character};
\end{tikzpicture}
\caption{A finite-dimensional representation of a torus decomposes
into weight spaces indexed by characters.}
\label{fig:torus-weight-space-decomposition}
\end{figure}

\begin{warningbox}
The theorem does not say that every element of \(V\) is a weight
vector. It says that every element is a sum of weight vectors. A vector
is itself a weight vector only when all of its nonzero components lie
in a single weight space.
\end{warningbox}

\subsection{Why the theorem is plausible}

We will not prove the weight-space decomposition here, but it is useful
to understand why one should expect it.

For \(T=\Gm^n\), the operators \(\rho(t)\) commute with one another
because \(T\) is abelian. More importantly, the algebraic nature of the
action forces the matrix coefficients to be built from characters.
The representation therefore separates into subspaces on which the
torus acts through individual Laurent monomials.

\begin{observation}
Another viewpoint uses the coordinate ring
\[
\Bbbk[T]=\Bbbk[M].
\]
The action gives \(V\) the structure of a comodule over
\(\Bbbk[M]\), and the basis elements \(\character{m}\) of the group
algebra separate the representation into the spaces \(V_m\). This is
the standard algebraic proof, but we will not need the comodule
language later.
\end{observation}

\begin{observation}
The weight-space decomposition holds for algebraic tori over an
algebraically closed field in arbitrary characteristic. Standard
references prove it as part of the representation theory of
diagonalizable algebraic groups.
\end{observation}

\subsection{Invariant vectors}

The zero element
\[
0\in M
\]
corresponds to the trivial character
\[
\character{0}=1.
\]
Its weight space is
\[
V_0
=
\{v\in V:t\cdot v=v\text{ for every }t\in T\}.
\]
This is precisely the subspace of \(T\)-invariant vectors, often
written
\[
V^T.
\]
Thus
\[
V^T=V_0.
\]

\begin{example}
For
\[
(t_1,t_2)\cdot(x_0,x_1,x_2)
=
(x_0,t_1x_1,t_2x_2),
\]
the invariant subspace is
\[
V^T=\Bbbk e_0.
\]
The other two coordinate lines have nonzero weights.
\end{example}

\begin{pauseandthink}
Let \(\Gm\) act on \(\Bbbk[x,y]_d\) by
\[
(t\cdot f)(x,y)=f(t^{-1}x,ty).
\]
Determine the weight of each monomial
\[
x^{d-i}y^i.
\]
For which values of \(d\) does the representation contain a nonzero
invariant vector?
\end{pauseandthink}

We have
\[
t\cdot x^{d-i}y^i
=
t^{-(d-i)}t^i x^{d-i}y^i
=
t^{2i-d}x^{d-i}y^i.
\]
Thus the weight is
\[
2i-d.
\]
A nonzero invariant monomial exists exactly when \(2i=d\), hence when
\(d\) is even.

\subsection{Equivariant linear maps preserve weights}

Let \(V\) and \(W\) be representations of \(T\).

\begin{definition}
A linear map
\[
F:V\longrightarrow W
\]
is \emph{\(T\)-equivariant} if
\[
F(t\cdot v)=t\cdot F(v)
\]
for every \(t\in T\) and \(v\in V\).
\end{definition}

If \(v\in V_m\), then
\[
\begin{aligned}
t\cdot F(v)
&=
F(t\cdot v)\\
&=
F(\character{m}(t)v)\\
&=
\character{m}(t)F(v).
\end{aligned}
\]
Therefore
\[
F(V_m)\subseteq W_m.
\]

So an equivariant map acts weight by weight. In particular, if \(W_m=0\),
then every equivariant map must vanish on \(V_m\).

\begin{example}
Suppose
\[
V=V_1\oplus V_3
\qquad\text{and}\qquad
W=W_1\oplus W_2
\]
are representations of \(\Gm\). A \(\Gm\)-equivariant map
\[
F:V\longrightarrow W
\]
may map \(V_1\) into \(W_1\), but it must vanish on \(V_3\), since
\(W_3=0\).
\end{example}

This simple observation is extremely useful. It turns many questions
about equivariant maps into separate linear-algebra problems on the
individual weight spaces.

\subsection{Subrepresentations and quotients}

A subspace \(U\subseteq V\) is a \emph{subrepresentation} if it is
preserved by \(T\). The weight decomposition restricts to
\[
U
=
\bigoplus_{m\in M}(U\cap V_m).
\]
Likewise, the quotient \(V/U\) inherits a representation, and its
weights are obtained from the corresponding quotients of weight
spaces.

Thus the representation theory of a torus has no hidden interaction
between distinct weights. Each weight may be studied independently.

\begin{observation}
This is one reason tori are among the simplest algebraic groups from
the viewpoint of representation theory. Their finite-dimensional
representations are controlled by:
\[
\text{a finite set of lattice points }m\in M
\quad\text{and}\quad
\text{their multiplicities }\dim V_m.
\]
\end{observation}

\subsection{Weights as a finite configuration in a lattice}

The set of weights
\[
\operatorname{Wt}(V)\subseteq M
\]
may be drawn as a finite configuration of lattice points in
\(M_{\RR}\). Multiplicity may be recorded by labels, repeated points,
or the dimensions of the associated weight spaces.

\begin{example}
Consider the representation of \(\Gm^2\) on \(\Bbbk^5\) with weights
\[
(0,0),\quad(1,0),\quad(1,0),\quad(0,1),\quad(2,-1).
\]
Then
\[
\operatorname{Wt}(V)
=
\{(0,0),(1,0),(0,1),(2,-1)\},
\]
with multiplicity \(2\) at \((1,0)\) and multiplicity \(1\) at each of
the other weights.
\end{example}

\begin{figure}[ht]
\centering
\begin{tikzpicture}[scale=1.15]
  \draw[->,black!55] (-0.8,0) -- (3.3,0);
  \draw[->,black!55] (0,-1.8) -- (0,2.2);

  \foreach \x in {0,1,2,3}
    \draw[black!20] (\x,-1.65) -- (\x,2.0);
  \foreach \y in {-1,0,1,2}
    \draw[black!20] (-0.65,\y) -- (3.1,\y);

  \fill[DictionaryBlue] (0,0) circle (2.6pt);
  \node[below left] at (0,0) {\((0,0)\)};

  \fill[PredictionTeal] (1,0) circle (3.4pt);
  \node[below] at (1,-0.12) {\((1,0)\)};
  \node[above] at (1,0.12) {mult.\ \(2\)};

  \fill[SurfaceGreen] (0,1) circle (2.6pt);
  \node[above left] at (0,1) {\((0,1)\)};

  \fill[ExampleGold] (2,-1) circle (2.6pt);
  \node[below right] at (2,-1) {\((2,-1)\)};

  \node[right] at (3.1,0) {\(M_{\RR}\)};
\end{tikzpicture}
\caption{A torus representation may be encoded by a finite set of
weights in \(M\), together with their multiplicities.}
\label{fig:finite-weight-configuration}
\end{figure}

This turns a representation into combinatorial data. Later, when a
torus acts on spaces of sections, the weights will often be the lattice
points of a polytope.

\subsection{Operations on representations}

The weight description makes standard constructions transparent.

If
\[
V=\bigoplus_m V_m
\qquad\text{and}\qquad
W=\bigoplus_{m'}W_{m'},
\]
then
\[
(V\oplus W)_m
=
V_m\oplus W_m.
\]

For tensor products,
\[
v\in V_m,\quad w\in W_{m'}
\quad\Longrightarrow\quad
v\otimes w\in(V\otimes W)_{m+m'}.
\]
Thus weights add under tensor product.

For dual representations,
\[
(V^\vee)_{-m}
\simeq
(V_m)^\vee.
\]
Thus weights change sign under duality.

\begin{example}
Let \(V\) and \(W\) be one-dimensional representations of weights
\(m\) and \(m'\), respectively. Then
\[
V\otimes W
\]
has weight
\[
m+m',
\]
while \(V^\vee\) has weight
\[
-m.
\]
This is the representation-theoretic version of
\[
\character{m}\character{m'}=\character{m+m'}
\qquad\text{and}\qquad
(\character{m})^{-1}=\character{-m}.
\]
\end{example}

\begin{prediction}
A torus acts on the space
\[
H^0(\PP^2,\cO_{\PP^2}(d))
\]
of homogeneous degree-\(d\) polynomials. What should its weight vectors
look like? How might the triangle \(d\Delta\cap M\) from Chapter~0
reappear in this representation?
\end{prediction}

\subsection{The basic classification}

The weight-space theorem may be summarized as a classification
statement.

\begin{theorem}[Classification of finite-dimensional torus
representations]
Up to isomorphism, a finite-dimensional representation of \(T\) is
determined by the multiplicities
\[
a_m=\dim V_m,
\qquad
m\in M,
\]
where \(a_m\) is a nonnegative integer and only finitely many \(a_m\)
are nonzero.

Equivalently,
\[
V
\simeq
\bigoplus_{m\in M}
\left(\Bbbk_m\right)^{\oplus a_m},
\]
where \(\Bbbk_m\) denotes the one-dimensional representation of
weight \(m\).
\end{theorem}

No extension data are needed: the representation is the direct sum of
its weight spaces.

\begin{observation}
Proofs of the weight-space theorem and this classification may be
found in standard treatments of diagonalizable algebraic groups or
linear algebraic groups; the Notes and references section lists the
sources used here.
\end{observation}

\begin{threedictionaries}
The representation theory of a torus is controlled by its character
lattice.

\medskip

\noindent
\begin{minipage}[t]{0.30\linewidth}
\centering
\textbf{Algebra}\par
\smallskip
A representation decomposes as
\[
V=\bigoplus_{m\in M}V_m.
\]
Equivariant maps preserve each weight.
\end{minipage}
\hfill
\begin{minipage}[t]{0.30\linewidth}
\centering
\textbf{Geometry}\par
\smallskip
After choosing a weight basis, the torus acts by diagonal matrices:
each coordinate direction is scaled independently.
\end{minipage}
\hfill
\begin{minipage}[t]{0.30\linewidth}
\centering
\textbf{Combinatorics}\par
\smallskip
A representation is encoded by finitely many lattice points
\[
m\in M
\]
and their multiplicities.
\end{minipage}
\end{threedictionaries}

\begin{guidingquestion}
How do the monomials, exponent vectors, one-parameter limits, and
lattice triangles of Chapter~0 fit into the character and weight
language developed in this chapter?
\end{guidingquestion}


\section*{What this chapter has accomplished}
\addcontentsline{toc}{section}{What this chapter has accomplished}
\markboth{What this chapter has accomplished}{What this chapter has accomplished}
\label{sec:why-torus-representation-matters}

Chapter~0 already contained nearly every object introduced in this
chapter. We first met them concretely in \(\PP^2\); we can now name
them and place them in the two lattices attached to its dense torus.

\subsection*{Pointers back to Chapter~0}

\begin{enumerate}[label=\textbf{\arabic*.},wide=0pt]

\item \textbf{Torus coordinates are bases of \(M\).}
On the chart \(U_0\), the functions
\[
z_1=\frac{x_1}{x_0},
\qquad
z_2=\frac{x_2}{x_0}
\]
are characters of the dense torus \(T\subseteq\PP^2\). Their lattice
points form a basis of
\[
M=X^*(T).
\]
The Laurent coordinate changes between \(U_0,U_1,U_2\) are changes of
basis in \(M\), hence are governed by
\(\operatorname{GL}(2,\ZZ)\).

\item \textbf{Exponent vectors are character lattice points.}
A Laurent monomial
\[
z_1^az_2^b
\]
is the character \(\character{m}\) corresponding, in the chosen basis,
to
\[
m=(a,b)\in M.
\]
Thus exponent vectors are coordinates for intrinsic characters.

\item \textbf{Monomial sections form weight bases.}
The degree-\(d\) monomials
\[
x_0^{d-a-b}x_1^ax_2^b,
\qquad
(a,b)\in d\Delta\cap M,
\]
form a weight basis of
\[
H^0\bigl(\PP^2,\cO_{\PP^2}(d)\bigr).
\]
Depending on the convention for the induced action on sections, the
weights may be the exponent points or their negatives. The essential
statement is unchanged: the lattice points of \(d\Delta\) index the
monomial weight spaces.

\item \textbf{Monomial curves are one-parameter subgroups.}
The parameterized curves
\[
t\longmapsto(t^{u_1},t^{u_2})
\]
are the one-parameter subgroups
\[
\oneps{u}:\Gm\longrightarrow T,
\qquad
u=(u_1,u_2)\in N=X_*(T).
\]
Characters live in \(M\); algebraic motions live in the dual lattice
\(N\).

\item \textbf{The pairing records orders along motions.}
For \(m\in M\) and \(u\in N\),
\[
\character{m}\bigl(\oneps{u}(t)\bigr)
=
t^{\pair{m}{u}}.
\]
Hence \(\pair{m}{u}\) is the order of vanishing or pole of the
character along the one-parameter motion. This is the bridge between
functions in \(M\) and directions in \(N\).

\i

tem \textbf{Fans organize boundary behavior.}
For
\[
t\longmapsto[1:t^{u_1}:t^{u_2}]
\]
in \(\PP^2\), vectors \(u\) lying in the same cone of the fan produce
compatible boundary limits. The rays
\[
(-1,-1),\qquad(1,0),\qquad(0,1)
\]
record the three torus-invariant boundary divisors, while the maximal
cones correspond to the three torus-fixed points.

The blowup of a fixed point inserts a new primitive ray. Thus the star
subdivision from Chapter~0 may now be read as adding a new
one-parameter direction, and geometrically as introducing the
exceptional divisor.

\end{enumerate}

\begin{threedictionaries}
The chapter may be summarized by one compact dictionary.

\medskip

\begin{center}
\small
\begin{tabularx}{\linewidth}{@{}
  >{\raggedright\arraybackslash}X
  >{\raggedright\arraybackslash}X
  >{\raggedright\arraybackslash}X@{}}
\toprule
\textbf{Algebra}
&
\textbf{Geometry}
&
\textbf{Combinatorics}
\\
\midrule
\(\Bbbk[T]=\Bbbk[M]\)
&
the torus \(T\)
&
the character lattice \(M\)
\\[1mm]
a character \(\character{m}\)
&
a monomial function or weight
&
a lattice point \(m\in M\)
\\[1mm]
a one-parameter subgroup \(\oneps{u}\)
&
a parameterized motion toward the boundary
&
a lattice point \(u\in N\)
\\[1mm]
\(\character{m}\circ\oneps{u}=t^{\pair{m}{u}}\)
&
order of vanishing or pole along the motion
&
the pairing \(M\times N\to\ZZ\)
\\
\bottomrule
\end{tabularx}
\end{center}
\end{threedictionaries}

\subsection*{The next construction}

The next chapter begins with the observation that a cone
\[
\sigma\subseteq N_{\RR}
\]
specifies a collection of allowed one-parameter directions. The
characters having no pole along any of them are precisely those
\(m\in M\) satisfying
\[
\pair{m}{u}\geq0
\qquad
\text{for every }u\in\sigma.
\]
They form the semigroup
\[
\sigma^\vee\cap M,
\]
and its semigroup algebra
\[
\Bbbk[\sigma^\vee\cap M]
\]
will produce an affine variety containing \(T\).

For the first quadrant in \(N_{\RR}\), this construction gives
\[
\sigma^\vee\cap M=\ZZ_{\geq0}^2,
\qquad
\Bbbk[\sigma^\vee\cap M]=\Bbbk[x,y],
\]
and therefore recovers \(\AAA^2\).

\begin{guidingquestion}
Given a cone of one-parameter directions in \(N_{\RR}\), what geometry
is encoded by the algebra \(\Bbbk[\sigma^\vee\cap M]\)?
\end{guidingquestion}


\clearpage

\begin{exercises}[Comprehension exercises]
\markboth{Comprehension exercises}{Comprehension exercises}

\item Determine all characters
\[
\chi:\Gm\longrightarrow\Gm.
\]
Then determine which regular maps of the form
\[
f(t)=c\,t^a
\qquad
(c\in\Bbbk^\times,\ a\in\ZZ)
\]
are characters. Explain the role of the identity element.

\item Let \(T=\Gm^2\), with coordinate characters \(x\) and \(y\).
For each of the following pairs, decide whether the two characters
form a basis of \(M=X^*(T)\). When they do, write the corresponding
invertible monomial change of coordinates.
\[
\begin{array}{lll}
\text{(a)}&x^2y,\ xy,&
\text{(b)}\ x^{-1},\ x^{-1}y,\\[1mm]
\text{(c)}&x^2,\ y,&
\text{(d)}\ xy,\ x^{-1}y.
\end{array}
\]

\item Consider the one-parameter subgroup
\[
\oneps{u}(t)
=
(t^{u_1},t^{u_2},t^{u_3})
\]
of \(\Gm^3\).
\begin{enumerate}[label=(\alph*),wide=0pt]
\item Determine when it extends to a morphism
\[
\AAA^1\longrightarrow\AAA^3.
\]
\item When it extends, describe its value at \(t=0\) in terms of the
vanishing pattern of \(u_1,u_2,u_3\).
\item Repeat the question for a translated curve
\[
t\longmapsto(a_1t^{u_1},a_2t^{u_2},a_3t^{u_3}),
\qquad
a_i\in\Bbbk^\times.
\]
\end{enumerate}

\item Let
\[
m=(3,-2,1)\in M\simeq\ZZ^3,
\qquad
u=(-1,4,2)\in N\simeq\ZZ^3.
\]
Compute \(\pair{m}{u}\), and interpret the answer as the order of
vanishing or pole of
\[
\character{m}\circ\oneps{u}
\]
at \(t=0\). Find a nonzero vector \(u'\in N\) satisfying
\[
\pair{m}{u'}=0.
\]

\item Let \(m_1,m_2\) be a basis of \(M\), with dual basis
\(u_1,u_2\) of \(N\). Define
\[
m'_1=2m_1+m_2,
\qquad
m'_2=m_1+m_2.
\]
\begin{enumerate}[label=(\alph*),wide=0pt]
\item Verify that \(m'_1,m'_2\) form a basis of \(M\).
\item Compute the basis \(u'_1,u'_2\) of \(N\) dual to
\(m'_1,m'_2\).
\item Check directly that
\[
\pair{m'_i}{u'_j}=\delta_{ij}.
\]
\end{enumerate}

\item Consider the torus morphism
\[
\phi:\Gm^2\longrightarrow\Gm^3,
\qquad
(t_1,t_2)
\longmapsto
(t_1^2t_2,\ t_1^{-1}t_2^3,\ t_1t_2^2).
\]
\begin{enumerate}[label=(\alph*),wide=0pt]
\item Write the matrix of
\[
\phi^*:M_{\Gm^3}\longrightarrow M_{\Gm^2}
\]
in the standard bases.
\item Write the matrix of
\[
\phi_*:N_{\Gm^2}\longrightarrow N_{\Gm^3}.
\]
\item Verify
\[
\pair{\phi^*(m)}{u}
=
\pair{m}{\phi_*(u)}
\]
for a general \(m\in\ZZ^3\) and \(u\in\ZZ^2\).
\end{enumerate}

\item Let \(T=\Gm^2\) act on \(V=\Bbbk^5\) by
\[
(t_1,t_2)\cdot
(v_1,v_2,v_3,v_4,v_5)
=
(t_1v_1,\ t_2v_2,\ t_1t_2v_3,\ t_1v_4,\ v_5).
\]
Determine every weight space, its weight, and its multiplicity.
Describe the invariant subspace \(V^T\).

\item Let
\[
V=V_{(1,0)}\oplus V_{(0,1)}\oplus V_{(1,1)}
\]
and
\[
W=W_{(1,0)}\oplus W_{(1,1)}
\]
be representations of \(T=\Gm^2\), with
\[
\dim V_{(1,0)}=2,\quad
\dim V_{(0,1)}=1,\quad
\dim V_{(1,1)}=3,
\]
and
\[
\dim W_{(1,0)}=1,\quad
\dim W_{(1,1)}=2.
\]
Compute the dimension of the vector space
\[
\Hom_T(V,W)
\]
of \(T\)-equivariant linear maps.

\item For the standard torus action on \(\PP^2\), list the monomial
basis of
\[
H^0\bigl(\PP^2,\cO_{\PP^2}(3)\bigr)
\]
and label each monomial by its exponent point in \(3\Delta\cap M\).
Using the standard pullback convention for the induced action on
sections, also record its weight. Which section spans the invariant
subspace?

\item Let
\[
\sigma
=
\Cone{(1,0),(1,2)}
\subseteq N_{\RR}\simeq\RR^2.
\]
Determine all
\[
m=(a,b)\in M\simeq\ZZ^2
\]
satisfying
\[
\pair{m}{u}\geq0
\qquad
\text{for every }u\in\sigma.
\]
Describe the semigroup \(\sigma^\vee\cap M\) by inequalities, and find
a finite collection of lattice points that generates it additively.

\end{exercises}

\clearpage
\section*{Challenge exercise}
\addcontentsline{toc}{section}{Challenge exercise}
\markboth{Challenge exercise}{Challenge exercise}

\noindent
\textbf{Projective limits, weights, and exposed faces.}

\medskip

Let \(T\) be an algebraic torus, and let
\[
V=\bigoplus_{m\in M}V_m
\]
be a finite-dimensional representation. Choose a nonzero vector
\[
v=\sum_{m\in S}v_m,
\qquad
0\neq v_m\in V_m,
\]
where \(S\subseteq M\) is the finite set of weights occurring in
\(v\). The torus acts on the projective space \(\PP(V)\).

Fix a one-parameter subgroup \(u\in N\). The goal is to determine
\[
\lim_{t\to0}\oneps{u}(t)\cdot[v]
\]
directly from the pairings
\[
\pair{m}{u},
\qquad
m\in S.
\]

\begin{enumerate}[label=\textbf{Hint \arabic*.},wide=0pt]

\item Show that
\[
\oneps{u}(t)\cdot v
=
\sum_{m\in S}t^{\pair{m}{u}}v_m.
\]

\item Let
\[
a=\min_{m\in S}\pair{m}{u}.
\]
Use projective rescaling by \(t^{-a}\) to show that the limit exists
and equals
\[
\left[
\sum_{\substack{m\in S\\
\pair{m}{u}=a}}
v_m
\right].
\]
Explain why this vector is nonzero.

\item Let
\[
P(v)=\Conv(S)\subseteq M_{\RR}
\]
be the convex hull of the weights occurring in \(v\). Show that the
weights contributing to the limit are precisely the points of \(S\)
lying on the face
\[
F_u
=
\left\{
m\in P(v):
\pair{m}{u}
=
\min_{m'\in P(v)}\pair{m'}{u}
\right\}.
\]
This is the face of \(P(v)\) exposed by the linear functional
\(m\mapsto\pair{m}{u}\).

\item Explain why replacing \(u\) by another vector \(u'\) that exposes
the same face produces a limit lying in the same projective weight
subspace
\[
\PP\left(\bigoplus_{m\in F_u\cap S}V_m\right).
\]

\item Specialize to
\[
V=\Bbbk^3
\]
with the \(\Gm^2\)-action
\[
(t_1,t_2)\cdot(x_0,x_1,x_2)
=
(x_0,t_1x_1,t_2x_2).
\]
The weights are
\[
0,\qquad m_1,\qquad m_2.
\]
For a point
\[
[x_0:x_1:x_2]\in\PP^2,
\]
use the formula above to recover the one-parameter limits described in
Chapter~0. Pay attention to the possibility that one or more
coordinates vanish, since the set \(S\) depends on the point.

\item Draw the weight triangle
\[
\Conv\{0,m_1,m_2\}\subseteq M_{\RR}.
\]
Partition \(N_{\RR}\) according to which face it exposes by taking the
minimum of \(\pair{m}{u}\). Compare the resulting decomposition
with the fan of \(\PP^2\).

\item Formulate a general conjecture: given a finite weight polytope
\(P\subseteq M_{\RR}\), how should the regions of vectors
\(u\in N_{\RR}\) exposing a fixed face of \(P\) fit together? What
familiar fan should they form?

\end{enumerate}

\clearpage
\section*{Notes and references}
\addcontentsline{toc}{section}{Notes and references}
\markboth{Notes and references}{Notes and references}

This chapter is intended as a working introduction to the language of
algebraic tori rather than a self-contained treatment of linear
algebraic groups. The classification of characters and
one-parameter subgroups, the perfect pairing between their lattices,
and the weight-space decomposition of a finite-dimensional torus
representation were stated with only brief indications of proof.

For the toric-geometric viewpoint on lattices, characters,
one-parameter subgroups, and group algebras, see the opening chapter
of Cox--Little--Schenck, \emph{Toric Varieties}. Standard references
for algebraic groups and the representation theory of tori include
Humphreys, \emph{Linear Algebraic Groups}; Borel,
\emph{Linear Algebraic Groups}; and Springer,
\emph{Linear Algebraic Groups}. In those texts, tori are treated as
examples of diagonalizable algebraic groups, and the weight-space
decomposition appears as part of the general representation theory of
diagonalizable groups.

Over a non-algebraically closed field, one must distinguish split
tori from tori that become split only after extending scalars. That
arithmetic direction lies outside the main path of these notes.
